\section{Mô hình hóa hệ thống}

\subsection{Danh mục ca sử dụng}\label{sec:system-modeling-overview}
Ba hình dưới đây không chỉ có vai trò minh họa mà còn là phần chuyển tiếp quan trọng trước khi đi vào đặc tả ca sử dụng. Chúng lần lượt trả lời ba câu hỏi nền tảng: \textbf{IRMS giao tiếp với ai}, \textbf{IRMS phải hỗ trợ những nhóm nghiệp vụ nào} và \textbf{một vòng vận hành đầu-cuối trong nhà hàng diễn ra ra sao}. Việc diễn giải ba hình này giúp phần ca sử dụng phía sau bám đúng bối cảnh toàn hệ thống thay vì bị nhìn như những kịch bản rời rạc; phần định vị và biện minh tương ứng được nối trực tiếp về Bảng~\ref{tab:appendix-traceability} và Bảng~\ref{tab:appendix-view-rationale} trong Phụ lục~\ref{appendix:traceability} và Phụ lục~\ref{appendix:view-rationale}.

\begin{figure}[H]
\centering
\includegraphics[width=\textwidth,height=0.84\textheight,keepaspectratio]{01_context_diagram_cua_irms.png}
\caption{Sơ đồ ngữ cảnh của IRMS}
\label{fig:irms-context-overview}
\end{figure}


Hình~\ref{fig:irms-context-overview} chốt rõ ranh giới hệ thống theo đúng kiểu \emph{context/process diagram}: IRMS được đặt ở trung tâm như một tiến trình điều phối duy nhất, còn xung quanh là các thực thể bên ngoài và hệ thống đối tác trao đổi dữ liệu với nó. Cách thể hiện này không chỉ trả lời câu hỏi \enquote{ai dùng IRMS}, mà còn làm rõ \enquote{dữ liệu gì đi vào, kết quả gì đi ra} ở từng điểm chạm như đặt bàn, xếp bàn, gọi món, điều phối bếp, quyết toán, kiểm soát truy cập và tích hợp dịch vụ ngoài.

\paragraph{Các điểm cần đọc từ sơ đồ ngữ cảnh.}
\begin{itemize}
  \item \textbf{Tuyến phục vụ phía trước}: \texttt{Customer}, \texttt{Host / Reception}, \texttt{Server / Floor Staff} và \texttt{Cashier} trao đổi trực tiếp các luồng dữ liệu thời gian thực như yêu cầu đặt bàn, trạng thái chờ, đơn gọi món, tách hóa đơn và nhập thông tin thanh toán.
  \item \textbf{Tuyến thực thi phía sau}: \texttt{Kitchen Operations} trả về trạng thái chế biến, tín hiệu món sẵn sàng hoặc hết nguyên liệu; \texttt{Manager} và \texttt{System Administrator} gửi thay đổi cấu hình, truy vấn báo cáo và yêu cầu kiểm toán ở lớp điều hành.
  \item \textbf{Các hợp đồng tích hợp ngoài}: \texttt{Payment Gateway}, \texttt{Notification Service} và \texttt{Receipt Printer} được đặt ngoài biên hệ thống để nhấn mạnh rằng thanh toán, gửi thông báo và in biên lai là các phụ thuộc cần được cách ly bằng giao tiếp có tên gọi rõ ràng.
  \item \textbf{Các luồng dữ liệu đều được đặt tên}: hình không dừng ở mức liệt kê vai trò, mà chỉ ra cụ thể các trao đổi như \enquote{reservation request}, \enquote{kitchen ticket}, \enquote{transaction result}, \enquote{audit trail} hoặc \enquote{print status}; nhờ đó ranh giới trách nhiệm của IRMS được khóa chặt hơn trước khi đi vào use case và kiến trúc.
\end{itemize}

Chính vì vậy, hình ngữ cảnh là cầu nối trực tiếp từ bối cảnh nghiệp vụ ở Chương 1 sang sơ đồ tổng quan kiến trúc ở Mục~\ref{sec:architecture-overview}: nó khóa ranh giới hệ thống trước khi tài liệu đi sâu vào cách phân rã các năng lực bên trong.

\begin{figure}[H]
\centering
\includegraphics[width=\textwidth,height=0.84\textheight,keepaspectratio]{02_use_case_overview_cua_irms.png}
\caption{Sơ đồ tổng quan ca sử dụng của IRMS}
\label{fig:irms-usecase-overview}
\end{figure}


Hình~\ref{fig:irms-usecase-overview} trình bày tổng quan ca sử dụng theo đúng ký pháp UML: các actor đứng ngoài biên hệ thống, các ca sử dụng được biểu diễn bằng ellipse và được gom thành năm gói nghiệp vụ lớn gồm đặt bàn và vòng đời bàn, gọi món và điều phối bếp, hóa đơn và thanh toán, tồn kho và dự báo nhập hàng, cùng quản trị, kiểm toán và phân tích. Cách tổ chức này cho thấy phạm vi của IRMS không dừng ở các thao tác trực tiếp trước khách hàng, mà bao trùm cả phần kết thúc vòng đời bàn, giao tiếp chủ động với khách hàng, điều phối ưu tiên trong bếp, hỗ trợ quyết định nhập hàng từ dữ liệu tiêu hao lịch sử và truy vết thao tác nhạy cảm ở tầng quản trị.

Nhìn theo góc độ nghiệp vụ, năm cụm này cũng tương ứng với năm lớp trách nhiệm khác nhau trong hệ thống. Cụm đặt bàn và bàn ăn quản lý năng lực tiếp nhận khách theo thời gian thực; cụm gọi món và bếp điều khiển luồng thực thi nóng nhất của nhà hàng; cụm hóa đơn và thanh toán chốt quyết định tài chính; cụm tồn kho và dự báo hỗ trợ điều hành vật tư; còn cụm quản trị, kiểm toán và phân tích bảo đảm hệ thống có thể được kiểm soát, truy vết và cải tiến sau vận hành. Vì vậy, sơ đồ tổng quan ca sử dụng không chỉ là danh mục chức năng, mà còn là bản đồ chỉ ra nơi nào cần phản hồi tức thời, nơi nào cần độ đúng đắn cao, và nơi nào có thể chấp nhận xử lý sau trên tuyến nền.

\paragraph{Năm gói ca sử dụng chính thể hiện trong hình.}
\begin{itemize}
  \item \textbf{Reservation and Seating}: gom toàn bộ vòng đời đặt bàn, check-in, danh sách chờ, giải phóng bàn và thông báo cho khách.
  \item \textbf{Ordering and Kitchen}: gom tuyến gọi món, tùy biến món, điều phối bếp, cập nhật trạng thái chế biến và xử lý ưu tiên trong giờ cao điểm.
  \item \textbf{Billing and Payment}: gom các thao tác tạo hóa đơn, tách hóa đơn, thanh toán, phát hành biên lai và hoàn tiền.
  \item \textbf{Inventory and Reorder}: gom tiêu hao tồn kho, cảnh báo mức tồn thấp và hỗ trợ quyết định nhập hàng.
  \item \textbf{Governance and Analytics}: gom các ca sử dụng quản trị menu, phân quyền, nhật ký kiểm toán, báo cáo và lịch ca.
\end{itemize}

Việc gom theo gói ở mức tổng quan giúp hình giữ được khả năng đọc, còn phần chi tiết của từng mã ca sử dụng vẫn được triển khai ngay sau đó trong danh mục và các đặc tả cụ thể. Cách trình bày này cũng tạo nhịp chuyển hợp lý sang các góc nhìn logic và mô-đun ở chương kiến trúc, nơi năm gói nghiệp vụ này được ánh xạ thành các miền trách nhiệm ổn định hơn; phần định vị và giải thích mở rộng của sơ đồ tổng quan ca sử dụng được dẫn ngay tới Bảng~\ref{tab:appendix-traceability} và Bảng~\ref{tab:appendix-view-rationale}.

\begin{figure}[H]
\centering
\includegraphics[width=\textwidth,height=0.84\textheight,keepaspectratio]{03_activity_diagram_end_to_end_restaurant_operation.png}
\caption{Sơ đồ hoạt động của luồng vận hành đầu-cuối nhà hàng}
\label{fig:irms-activity-end-to-end}
\end{figure}


Hình~\ref{fig:irms-activity-end-to-end} mô tả luồng vận hành đầu-cuối tiêu biểu của nhà hàng, bắt đầu từ khâu tiếp nhận khách hoặc kiểm tra đặt bàn, chuyển sang mở phiên bàn, tạo đơn gọi món, gửi phiếu xuống bếp, cập nhật tiến độ chế biến, giao món, lập hóa đơn, xử lý thanh toán và kết thúc bằng giải phóng bàn. Điểm quan trọng của sơ đồ hoạt động này là nó cho thấy những trạng thái nào phải được đồng bộ mạnh giữa các bộ phận: một bàn chỉ được mở một phiên phục vụ, một đơn đã xác nhận phải sinh phiếu bếp, và một hóa đơn chỉ được đóng khi thanh toán hoàn tất hoặc được ghi nhận là thất bại có kiểm soát.

Nhìn dưới góc độ kiến trúc, luồng đầu-cuối này cũng chỉ ra ranh giới giữa \textbf{đường giao dịch nóng} và \textbf{đường xử lý nền}. Các bước như xác nhận đơn gọi món, cập nhật trạng thái phiếu \texttt{Ready}, ghi nhận thanh toán thành công và đóng \texttt{TableSession} thuộc về tuyến đồng bộ bắt buộc vì chúng tác động trực tiếp đến trải nghiệm khách hàng. Ngược lại, các bước như cập nhật báo cáo, phát cảnh báo mức tồn thấp, gửi thông báo xác nhận hoặc đồng bộ số liệu quản trị có thể được đẩy sang luồng bất đồng bộ để giảm tải cho các điểm chạm thời gian thực.

\paragraph{Những mắt xích kiến trúc được khóa bởi sơ đồ đầu-cuối.}
\begin{itemize}
  \item \textbf{Điểm bàn giao giữa các khu vực}: lễ tân mở phiên bàn, phục vụ xác nhận đơn, bếp cập nhật trạng thái, thu ngân quyết toán; mỗi điểm bàn giao đều gắn với một thay đổi trạng thái có hậu quả nghiệp vụ.
  \item \textbf{Điểm đồng bộ bắt buộc}: mở \texttt{TableSession}, xác nhận đơn, đánh dấu món sẵn sàng, ghi nhận thanh toán và giải phóng bàn không được nhập nhằng về trạng thái.
  \item \textbf{Điểm bất đồng bộ chấp nhận được}: thông báo, báo cáo và một phần cảnh báo tồn kho được đẩy ra sau mà không làm sai tuyến phục vụ trước mặt khách hàng.
  \item \textbf{Điểm nối sang chương kiến trúc}: hình này là cơ sở để kiểm tra xem các sơ đồ logic, mô-đun, thành phần và triển khai ở chương sau có thật sự bảo vệ được luồng vận hành đầu-cuối hay không.
\end{itemize}

\begin{table}[H]
\centering
\caption{Danh mục 25 ca sử dụng chính được đặc tả chi tiết}
\begin{tabularx}{\textwidth}{>{\bfseries}p{4.0cm} X}
\toprule
Nhóm & Mã ca sử dụng \\
\midrule
Nhóm đặt bàn và bàn ăn & UC-RES-01, UC-RES-02, UC-RES-03, UC-RES-04, UC-RES-05 \\
Nhóm gọi món và điều phối bếp & UC-ORD-01, UC-ORD-02, UC-ORD-03, UC-KIT-01, UC-KIT-02, UC-KIT-03, UC-ORD-04 \\
Nhóm thanh toán & UC-BIL-01, UC-BIL-02, UC-PAY-01, UC-PAY-02, UC-PAY-03 \\
Nhóm tồn kho, cảnh báo và dự báo & UC-INV-01, UC-INV-02, UC-INV-03 \\
Nhóm quản trị, kiểm toán và phân tích & UC-REP-01, UC-ADM-01, UC-ADM-02, UC-ADM-03, UC-OPS-01 \\
\bottomrule
\end{tabularx}
\end{table}

Năm ca sử dụng giữ vai trò khóa kín phạm vi nghiệp vụ ở phần này là \texttt{UC-RES-04}, \texttt{UC-RES-05}, \texttt{UC-KIT-03}, \texttt{UC-INV-03} và \texttt{UC-ADM-03}. Chúng làm rõ các mắt xích dễ bị xem nhẹ khi mô tả hệ thống nhà hàng ở mức khái quát: kết thúc vòng đời bàn sau thanh toán, giao tiếp chủ động với khách hàng, điều phối ưu tiên trong bếp, hỗ trợ quyết định nhập hàng từ dữ liệu lịch sử và truy vết thao tác nhạy cảm ở tầng quản trị.

\subsection{Góc nhìn sử dụng}
Góc nhìn sử dụng cho thấy cách các tác nhân tương tác với các dịch vụ qua các kịch bản nghiệp vụ tiêu biểu.
Về mặt lý thuyết, góc nhìn sử dụng là công cụ rất tốt để \textbf{kiểm chứng kiến trúc bằng kịch bản}.
Một phương án phân rã đẹp trên góc nhìn mô-đun vẫn có thể thất bại khi đi qua tình huống thực,
nếu chuỗi tương tác quá dài, lỗi khó giải thích hoặc ranh giới giao dịch không rõ.
Do đó, bốn kịch bản dưới đây được chọn vì chúng chạm vào bốn luồng cốt lõi nhất của IRMS.

\subsubsection{Kịch bản 1: Từ đặt bàn đến tiếp nhận và xếp bàn}
\begin{enumerate}
  \item Khách hàng hoặc lễ tân tạo đặt bàn.
  \item \texttt{Reservation Service} kiểm tra khả dụng và gán bàn hoặc giữ chỗ.
  \item Khi khách đến, \texttt{Host UI} xác nhận tiếp nhận và mở \texttt{TableSession}.
  \item \texttt{TableSession} được chuyển tới \texttt{Server POS UI} để phục vụ gọi món.
\end{enumerate}

\subsubsection{Kịch bản 2: Từ gọi món đến điều phối bếp và phục vụ}
\begin{enumerate}
  \item Phục vụ tạo đơn gọi món và cấu hình tùy chọn món hoặc gợi ý dị ứng trên \texttt{Server POS UI}.
  \item \texttt{Ordering Service} xác nhận đơn gọi món và phát sự kiện tạo phiếu bếp.
  \item \texttt{Kitchen Coordination Service} tách phiếu theo trạm bếp và hiển thị lên \texttt{Kitchen Display UI}.
  \item Bếp cập nhật trạng thái từng dòng phiếu; màn hình gọi món nhận lại trạng thái sẵn sàng.
  \item Phục vụ giao món và đánh dấu đã phục vụ; \texttt{Inventory Service} nhận sự kiện để trừ tồn.
\end{enumerate}

\subsubsection{Kịch bản 3: Từ hóa đơn đến thanh toán và biên lai}
\begin{enumerate}
  \item Thu ngân hoặc phục vụ gom các đơn gọi món của bàn thành hóa đơn.
  \item \texttt{Billing Service} áp dụng khuyến mãi, tiền boa và tách hóa đơn nếu cần.
  \item Thanh toán được thực hiện qua tiền mặt hoặc cổng thanh toán điện tử.
  \item Khi thanh toán đủ, biên lai được phát hành và \texttt{TableSession} chuyển sang chuẩn bị đóng.
\end{enumerate}

\subsubsection{Kịch bản 4: Từ cảnh báo mức tồn thấp đến quyết định của quản lý}
\begin{enumerate}
  \item \texttt{Inventory Service} cập nhật lượng tồn thực tế sau mỗi món đã phục vụ hoặc sau mỗi lần điều chỉnh thủ công.
  \item Nếu xuống dưới ngưỡng, hệ thống tạo sự kiện cảnh báo mức tồn thấp.
  \item \texttt{Manager Portal} nhận cảnh báo, hiển thị nguyên liệu bị ảnh hưởng và món có nguy cơ hết hàng.
  \item Quản lý quyết định nhập thêm hàng hoặc tạm ẩn món trong menu.
\end{enumerate}

\begin{infobox}
\textbf{Giải thích chi tiết cho góc nhìn sử dụng.}
Góc nhìn sử dụng đặc biệt quan trọng với IRMS vì hệ thống không phải một khối xử lý tĩnh; nó là một chuỗi phối hợp giữa người và máy theo thời gian thực. Mỗi kịch bản ở trên đều cho thấy ranh giới nào cần đồng bộ mạnh, chẳng hạn như thanh toán, và ranh giới nào chấp nhận xử lý bất đồng bộ, chẳng hạn như báo cáo hoặc cảnh báo mức tồn thấp.
\end{infobox}

\begin{landscape}
\begin{figure}[H]
\centering
\IfFileExists{11_sequence_diagram_reservation_check_in_and_seating.pdf}{
  \includegraphics[width=0.97\linewidth,height=0.84\textheight,keepaspectratio]{11_sequence_diagram_reservation_check_in_and_seating.pdf}
}{
  \IfFileExists{11_sequence_diagram_reservation_check_in_and_seating.png}{
    \includegraphics[width=0.97\linewidth,height=0.84\textheight,keepaspectratio]{11_sequence_diagram_reservation_check_in_and_seating.png}
  }{
    \fbox{\parbox{0.82\linewidth}{Thiếu file \texttt{11\_sequence\_diagram\_reservation\_check\_in\_and\_seating.(pdf/png)}.}}
  }
}
\caption{Sơ đồ trình tự cho kịch bản từ đặt bàn đến tiếp nhận và xếp bàn}
\label{fig:irms-usage-reservation}
\end{figure}
\end{landscape}


Hình~\ref{fig:irms-usage-reservation} cho thấy \texttt{Reservation} và \texttt{TableSession}
không phải là cùng một thực thể nghiệp vụ dù chúng liên quan trực tiếp.
Điểm mạnh của cách tách này là hệ thống xử lý được nhiều trường hợp thực tế:
khách có đặt bàn nhưng không đến,
khách vào trực tiếp không có đặt bàn nhưng vẫn mở phiên phục vụ,
hoặc lễ tân phải đổi bàn khi năng lực phục vụ thay đổi.
Nếu trộn hai khái niệm này thành một,
kiến trúc sẽ khó biểu diễn các trạng thái chuyển tiếp quan trọng của vận hành nhà hàng.

Đánh đổi là luồng xếp bàn phải cẩn thận hơn về đồng bộ trạng thái:
đặt bàn có thể còn hiệu lực nhưng bàn ăn chưa sẵn sàng;
bàn có thể vừa được giải phóng nhưng vẫn đang dọn dẹp;
thao tác tiếp nhận có thể đến muộn so với thời điểm giữ chỗ ban đầu.
Do đó, góc nhìn sử dụng này nhắc rằng đặt bàn và xếp bàn tuy nhìn đơn giản,
nhưng thực ra là nơi chứa nhiều quy tắc nghiệp vụ về thời gian và năng lực bàn.

\paragraph{Các thành phần chính thể hiện trong hình.}
\begin{itemize}
  \item \textbf{Khách hàng và Host UI}: khởi tạo yêu cầu đặt bàn hoặc check-in.
  \item \textbf{Reservation Service}: chịu trách nhiệm xác thực trạng thái đặt bàn và gợi ý bàn.
  \item \textbf{TableSession và Server POS UI}: nhận bàn đã mở để chuyển sang luồng phục vụ thực tế.
\end{itemize}

\paragraph{Lý do thiết kế của sơ đồ này.}
\begin{itemize}
  \item Sơ đồ làm rõ ranh giới giữa dữ liệu đặt bàn và phiên bàn thực tế.
  \item Đây là kịch bản giúp chứng minh tại sao bài toán xếp bàn cần được mô hình hóa riêng thay vì gộp vào một trạng thái bàn đơn giản.
\end{itemize}

\begin{landscape}
\begin{figure}[H]
\centering
\IfFileExists{12_sequence_diagram_place_order_to_kitchen.pdf}{
  \includegraphics[width=0.97\linewidth,height=0.84\textheight,keepaspectratio]{12_sequence_diagram_place_order_to_kitchen.pdf}
}{
  \IfFileExists{12_sequence_diagram_place_order_to_kitchen.png}{
    \includegraphics[width=0.97\linewidth,height=0.84\textheight,keepaspectratio]{12_sequence_diagram_place_order_to_kitchen.png}
  }{
    \fbox{\parbox{0.82\linewidth}{Thiếu file \texttt{12\_sequence\_diagram\_place\_order\_to\_kitchen.(pdf/png)}.}}
  }
}
\caption{Sơ đồ trình tự cho kịch bản từ gọi món đến điều phối bếp và phục vụ}
\label{fig:irms-usage-order-kitchen}
\end{figure}
\end{landscape}


Hình~\ref{fig:irms-usage-order-kitchen} là kịch bản quan trọng nhất để kiểm chứng \textbf{khả năng đáp ứng}.
Nguyên tắc thiết kế nên là:
\emph{lưu đơn gọi món thành công trước, rồi mới lan truyền sang bếp}.
Điều này bảo vệ tuyến đầu theo hai cách.
Thứ nhất, màn hình gọi món nhận phản hồi nhanh hơn vì không phải đợi toàn bộ chuỗi bếp, tồn kho và báo cáo.
Thứ hai, khi việc chuyển sang bếp gặp trục trặc,
hệ thống vẫn có thể thử lại từ trạng thái đơn đã xác nhận thay vì mất dữ liệu đơn gọi món.

Ưu điểm của mô hình này là tách được
\textbf{bước chốt giao dịch} khỏi \textbf{bước lan truyền vận hành}.
Tuy nhiên, đánh đổi là màn hình bếp và màn hình gọi món phải xử lý được trạng thái trung gian:
đơn đã xác nhận nhưng phiếu chưa hiện lên màn hình bếp ngay lập tức;
một dòng món có thể đã sẵn sàng ở bếp nhưng màn hình gọi món chưa nhận được cập nhật cuối cùng;
việc trừ tồn cũng có thể đi sau sự kiện phục vụ xong.
Vì thế, góc nhìn sử dụng này là bằng chứng rằng cơ chế nhất quán trễ ở IRMS phải được dùng có kiểm soát, không dùng tùy tiện.

\paragraph{Các thành phần chính thể hiện trong hình.}
\begin{itemize}
  \item \textbf{Server POS UI}: là điểm bắt đầu của thao tác gọi món.
  \item \textbf{Ordering Service}: là nơi chốt giao dịch gọi món.
  \item \textbf{Kitchen Coordination Service và Kitchen Display UI}: nhận luồng lan truyền sang bếp sau khi đơn đã được xác nhận.
  \item \textbf{Inventory Service}: chỉ nhận cập nhật sau khi món đã đi qua luồng phục vụ phù hợp.
\end{itemize}

\paragraph{Lý do thiết kế của sơ đồ này.}
\begin{itemize}
  \item Sơ đồ này nhấn mạnh thứ tự \enquote{xác nhận trước, lan truyền sau} của tuyến gọi món.
  \item Đây là kịch bản quan trọng nhất để giải thích vì sao kiến trúc phải cân bằng giữa phản hồi nhanh và nhất quán có kiểm soát.
\end{itemize}

\begin{landscape}
\begin{figure}[H]
\centering
\IfFileExists{13_sequence_diagram_checkout_and_payment.pdf}{
  \includegraphics[width=0.97\linewidth,height=0.84\textheight,keepaspectratio]{13_sequence_diagram_checkout_and_payment.pdf}
}{
  \IfFileExists{13_sequence_diagram_checkout_and_payment.png}{
    \includegraphics[width=0.97\linewidth,height=0.84\textheight,keepaspectratio]{13_sequence_diagram_checkout_and_payment.png}
  }{
    \fbox{\parbox{0.82\linewidth}{Thiếu file \texttt{13\_sequence\_diagram\_checkout\_and\_payment.(pdf/png)}.}}
  }
}
\caption{Sơ đồ trình tự cho kịch bản từ hóa đơn đến thanh toán và biên lai}
\label{fig:irms-usage-payment}
\end{figure}
\end{landscape}


Hình~\ref{fig:irms-usage-payment} là nơi kiến trúc phải nghiêng mạnh về \textbf{độ tin cậy}.
Ở kịch bản này, khối lập hóa đơn không chỉ làm phép cộng cuối cùng,
mà còn là nơi áp dụng khuyến mãi, tiền boa, tách hóa đơn, chọn phương thức thanh toán,
giao tiếp với cổng thanh toán, phát hành biên lai và đổi trạng thái phiên phục vụ.
Đây là luồng có mật độ chuyển trạng thái cao nhất và hệ quả nghiệp vụ lớn nhất nếu xử lý sai.

Ưu điểm của việc giữ thanh toán trong luồng đồng bộ là hệ thống có thể trả lời rất rõ:
giao dịch thành công hay thất bại, hóa đơn còn mở hay đã đóng, phiên phục vụ còn tiếp diễn hay đã kết thúc.
Ngược lại, nếu thanh toán bị đẩy sang bất đồng bộ quá mức,
thu ngân và khách hàng sẽ rất khó hiểu trạng thái thật sự của giao dịch.
Nhược điểm của cách làm đồng bộ là đường xử lý thanh toán trở thành điểm nhạy cảm với hệ ngoài,
đặc biệt khi cổng thanh toán chậm hoặc lỗi.
Vì thế, đây là nơi cần thời gian chờ rõ ràng, khóa chống xử lý lặp, nhật ký kiểm toán và cơ chế xử lý lỗi minh bạch nhất.

\paragraph{Các thành phần chính thể hiện trong hình.}
\begin{itemize}
  \item \textbf{Cashier UI}: là điểm điều khiển nghiệp vụ của thu ngân.
  \item \textbf{Billing Service}: giữ trung tâm của tuyến quyết toán.
  \item \textbf{Payment Gateway Adapter}: đại diện cho tích hợp thanh toán ngoài hệ thống.
  \item \textbf{Receipt Printer Adapter và TableSession}: kết thúc luồng bằng phát hành biên lai và đóng phiên bàn.
\end{itemize}

\paragraph{Lý do thiết kế của sơ đồ này.}
\begin{itemize}
  \item Đây là luồng nhạy cảm nhất về độ đúng đắn nghiệp vụ, nên cần được tách riêng để làm rõ ranh giới trách nhiệm và ranh giới giao dịch.
  \item Việc tách rõ cổng thanh toán, phát hành biên lai và đóng phiên bàn giúp người đọc thấy ranh giới trách nhiệm của từng bước sau khi thanh toán hoàn tất.
\end{itemize}

\begin{landscape}
\begin{figure}[H]
\centering
\IfFileExists{14_sequence_diagram_stock_update_and_low_stock_alert.pdf}{
  \includegraphics[width=0.97\linewidth,height=0.84\textheight,keepaspectratio]{14_sequence_diagram_stock_update_and_low_stock_alert.pdf}
}{
  \IfFileExists{14_sequence_diagram_stock_update_and_low_stock_alert.png}{
    \includegraphics[width=0.97\linewidth,height=0.84\textheight,keepaspectratio]{14_sequence_diagram_stock_update_and_low_stock_alert.png}
  }{
    \fbox{\parbox{0.82\linewidth}{Thiếu file \texttt{14\_sequence\_diagram\_stock\_update\_and\_low\_stock\_alert.(pdf/png)}.}}
  }
}
\caption{Sơ đồ trình tự cho kịch bản từ cảnh báo mức tồn thấp đến quyết định của quản lý}
\label{fig:irms-usage-stock}
\end{figure}
\end{landscape}


Hình~\ref{fig:irms-usage-stock} cho thấy rõ vì sao quản lý tồn kho và gửi thông báo không nên chen trực tiếp vào tuyến giao dịch nóng.
Việc theo dõi tồn kho cần chính xác ở mức vận hành,
nhưng không nhất thiết phải khóa trải nghiệm gọi món của phục vụ
chỉ vì một cảnh báo quản lý chưa được gửi xong.
Cách đặt việc trừ tồn và cảnh báo mức tồn thấp trên luồng sự kiện
giúp hệ thống giữ được \textbf{khả năng sửa đổi}:
khi chính sách trừ kho, ánh xạ công thức hay ngưỡng cảnh báo thay đổi,
phần ảnh hưởng chủ yếu nằm trong khối tồn kho và khối gửi thông báo.

Nhược điểm là lượng tồn hiển thị cho quản lý có thể hơi chậm so với giao dịch mới nhất
nếu bộ xử lý nền bị trễ hoặc sự kiện đang được thử lại.
Điều này chấp nhận được trong phạm vi IRMS,
miễn là không nhầm lẫn cảnh báo tồn kho với xác nhận thanh toán.
Nói cách khác, góc nhìn sử dụng này giúp chứng minh nơi nào được phép chấp nhận độ trễ đồng bộ,
và nơi nào thì tuyệt đối không nên chấp nhận điều đó.

\paragraph{Các thành phần chính thể hiện trong hình.}
\begin{itemize}
  \item \textbf{Ordering Service}: là nguồn phát sinh biến động vận hành cần suy diễn sang kho.
  \item \textbf{Inventory Service và LowStockRule}: chịu trách nhiệm cập nhật tồn và đánh giá ngưỡng cảnh báo.
  \item \textbf{Notification Adapter và Manager UI}: là đầu ra quản trị của luồng cảnh báo.
\end{itemize}

\paragraph{Lý do thiết kế của sơ đồ này.}
\begin{itemize}
  \item Sơ đồ cho thấy rõ tồn kho và cảnh báo nên đi trên một tuyến hỗ trợ thay vì chen vào đường giao dịch nóng.
  \item Đây là ví dụ rõ nhất để giải thích sự khác nhau giữa \textbf{đúng về vận hành} và \textbf{phải phản hồi tức thời} trong kiến trúc.
\end{itemize}

\subsection{Đặc tả ca sử dụng}


\usepackage{longtable}
\usepackage{enumitem}

\renewcommand{\arraystretch}{1.3}

\begin{longtable}{|p{4cm}|p{10cm}|}
\hline
\textbf{Use Case ID} & UC-RES-01 \\ \hline
\textbf{Use Case Name} & Tạo đặt bàn \\ \hline
\textbf{Primary Actors} & Khách hàng, Lễ tân \\ \hline

\textbf{Description} & Tạo một đặt bàn hợp lệ cho ngày, giờ và số lượng khách cụ thể. \\ \hline

\textbf{Trigger} & Khách hàng gửi yêu cầu đặt bàn qua quầy lễ tân, điện thoại hoặc web/app. \\ \hline

\textbf{Preconditions} &
\begin{itemize}[leftmargin=*, noitemsep]
  \item Nhà hàng có mở nhận đặt bàn trong khung giờ yêu cầu.
  \item Có tối thiểu tên và số điện thoại khách hàng.
  \item Sơ đồ bàn và năng lực phục vụ đã được cấu hình.
\end{itemize} \\ \hline

\textbf{Postconditions} &
\begin{itemize}[leftmargin=*, noitemsep]
  \item Tạo đặt bàn với trạng thái \texttt{Pending} hoặc \texttt{Confirmed}.
  \item Hệ thống giữ chỗ hoặc gợi ý bàn phù hợp.
  \item Khách nhận mã đặt bàn và thông báo xác nhận.
\end{itemize} \\ \hline

\textbf{Normal Flow} &
\begin{enumerate}[leftmargin=*, noitemsep]
  \item Lễ tân nhập thời gian, số khách, tên và số điện thoại.
  \item Hệ thống kiểm tra và validate dữ liệu đầu vào.
  \item Hệ thống kiểm tra sức chứa, bàn trống và xung đột đặt bàn.
  \item Hệ thống gợi ý bàn hoặc nhóm bàn phù hợp.
  \item Lễ tân chọn bàn và nhập ghi chú nếu có.
  \item Hệ thống tạo đặt bàn và sinh mã đặt bàn.
  \item Hệ thống gửi thông báo xác nhận cho khách.
\end{enumerate} \\ \hline

\textbf{Alternative Flows} &
\begin{minipage}[t]{\linewidth}
\textbf{A1. Không đủ sức chứa (từ bước 3)} \\
3.1 Hệ thống đề xuất khung giờ khác hoặc đưa khách vào danh sách chờ.

\vspace{0.5em}
\textbf{A2. Yêu cầu đặc biệt (từ bước 4)} \\
4.1 Lễ tân lọc danh sách bàn theo khu vực hoặc thuộc tính yêu cầu.

\vspace{0.5em}
\textbf{A3. Không gán được bàn (từ bước 5)} \\
5.1 Hệ thống tạo đặt bàn trạng thái \texttt{Pending} và giữ chỗ theo số lượng khách.
\end{minipage} \\ \hline

\textbf{Exception Flows} &
\begin{minipage}[t]{\linewidth}
\textbf{E1. Dữ liệu không hợp lệ (từ bước 2)} \\
2.1 Hệ thống yêu cầu nhập lại thông tin khách.

\vspace{0.5em}
\textbf{E2. Lỗi đồng bộ sơ đồ bàn (từ bước 3)} \\
3.2 Hệ thống khóa thao tác và yêu cầu tải lại dữ liệu.
\end{minipage} \\ \hline

\textbf{Business Rules} &
\begin{itemize}[leftmargin=*, noitemsep]
  \item Không vượt quá sức chứa tối đa của bàn hoặc nhóm bàn.
  \item Nhóm khách lớn phải được xác nhận bởi lễ tân hoặc quản lý.
\end{itemize} \\ \hline

\end{longtable}



   \renewcommand{\arraystretch}{1.3}

\begin{longtable}{|p{4cm}|p{10cm}|}
\hline
\textbf{Use Case ID} & UC-RES-02 \\ \hline
\textbf{Use Case Name} & Nhận khách và sắp bàn \\ \hline
\textbf{Primary Actors} & Lễ tân, Phục vụ \\ \hline

\textbf{Description} & Tiếp nhận khách đến nhà hàng và gán khách vào bàn hoặc nhóm bàn phù hợp. \\ \hline

\textbf{Trigger} & Khách đến quầy lễ tân hoặc được gọi từ danh sách chờ. \\ \hline

\textbf{Preconditions} &
\begin{itemize}[leftmargin=*, noitemsep]
  \item Khách có đặt bàn hoặc có thể được xếp chỗ tại chỗ.
  \item Sơ đồ bàn và trạng thái bàn được đồng bộ theo thời gian thực.
\end{itemize} \\ \hline

\textbf{Postconditions} &
\begin{itemize}[leftmargin=*, noitemsep]
  \item Một phiên phục vụ bàn (\texttt{TableSession}) được mở.
  \item Bàn được chuyển sang trạng thái \texttt{Occupied} hoặc \texttt{Reserved-Arrived}.
  \item Phục vụ nhận được thông báo tiếp nhận bàn mới.
\end{itemize} \\ \hline

\textbf{Normal Flow} &
\begin{enumerate}[leftmargin=*, noitemsep]
  \item Lễ tân tìm đặt bàn theo mã, số điện thoại hoặc tên khách.
  \item Hệ thống hiển thị thông tin đặt bàn, trạng thái, thời điểm đến và bàn gợi ý.
  \item Lễ tân xác nhận khách đến, số khách thực tế và thời gian check-in.
  \item Hệ thống kiểm tra hiệu lực đặt bàn và trạng thái bàn.
  \item Lễ tân điều chỉnh hoặc ghép bàn nếu số lượng khách thay đổi.
  \item Hệ thống mở phiên phục vụ (\texttt{TableSession}).
  \item Hệ thống cập nhật trạng thái bàn thành \texttt{Occupied} hoặc \texttt{Reserved-Arrived}.
  \item Hệ thống phân công bàn cho phục vụ theo khu vực.
  \item Phục vụ nhận thông báo và chuẩn bị phục vụ.
\end{enumerate} \\ \hline

\textbf{Alternative Flows} &
\begin{minipage}[t]{\linewidth}

\textbf{A1. Khách vãng lai (từ bước 1)} \\
1.1 Lễ tân chọn chức năng xếp bàn trực tiếp. \\
1.2 Hệ thống gợi ý bàn trống phù hợp. \\
1.3 Lễ tân chọn bàn và tiếp tục từ bước 5.

\vspace{0.5em}
\textbf{A2. Thay đổi số lượng khách (từ bước 3)} \\
3.1 Lễ tân cập nhật số khách thực tế. \\
3.2 Hệ thống gợi ý lại bàn phù hợp hoặc yêu cầu ghép bàn.

\end{minipage} \\ \hline

\textbf{Exception Flows} &
\begin{minipage}[t]{\linewidth}

\textbf{E1. Bàn chưa sẵn sàng (từ bước 4)} \\
4.1 Hệ thống đề xuất bàn thay thế hoặc đưa khách vào danh sách chờ.

\vspace{0.5em}
\textbf{E2. Đặt bàn hết hiệu lực (từ bước 4)} \\
4.2 Hệ thống yêu cầu xác nhận phục hồi từ lễ tân hoặc quản lý.

\vspace{0.5em}
\textbf{E3. Lỗi đồng bộ sơ đồ bàn (từ bước 4)} \\
4.3 Hệ thống tạm dừng thao tác và yêu cầu tải lại dữ liệu.

\end{minipage} \\ \hline

\textbf{Business Rules} &
\begin{itemize}[leftmargin=*, noitemsep]
  \item Mỗi bàn chỉ có một \texttt{TableSession} hoạt động tại một thời điểm.
  \item Không được gán bàn nếu bàn đang được sử dụng.
  \item Thời gian giữ chỗ tối đa được cấu hình theo ca phục vụ.
\end{itemize} \\ \hline

\end{longtable}



    \renewcommand{\arraystretch}{1.3}

\begin{longtable}{|p{4cm}|p{10cm}|}
\hline
\textbf{Use Case ID} & UC-RES-03 \\ \hline
\textbf{Use Case Name} & Quản lý danh sách chờ \\ \hline
\textbf{Primary Actors} & Lễ tân \\ \hline

\textbf{Description} & Ghi nhận khách chờ bàn và điều phối khách khi có bàn phù hợp. \\ \hline

\textbf{Trigger} & Nhà hàng hết bàn phù hợp hoặc khách bị trì hoãn. \\ \hline

\textbf{Preconditions} &
\begin{itemize}[leftmargin=*, noitemsep]
  \item Lễ tân có quyền truy cập danh sách chờ.
  \item Hệ thống có khả năng ước tính thời gian chờ.
\end{itemize} \\ \hline

\textbf{Postconditions} &
\begin{itemize}[leftmargin=*, noitemsep]
  \item Khách được thêm vào danh sách chờ với trạng thái \texttt{Waiting}.
  \item Khách được gọi theo thứ tự ưu tiên khi có bàn phù hợp.
\end{itemize} \\ \hline

\textbf{Normal Flow} &
\begin{enumerate}[leftmargin=*, noitemsep]
  \item Lễ tân nhập thông tin khách, số lượng khách và kênh liên hệ.
  \item Hệ thống ước tính thời gian chờ dựa trên tình trạng bàn.
  \item Lễ tân xác nhận thêm khách vào danh sách chờ.
  \item Khi có bàn trống, hệ thống tìm khách phù hợp theo sức chứa và ưu tiên.
  \item Lễ tân chọn khách để gọi.
  \item Hệ thống gửi thông báo và cập nhật trạng thái \texttt{Called}.
  \item Lễ tân xác nhận khách đến và chuyển sang use case Nhận khách và sắp bàn.
\end{enumerate} \\ \hline

\textbf{Alternative Flows} &
\begin{minipage}[t]{\linewidth}

\textbf{A1. Ưu tiên khách (từ bước 4)} \\
4.1 Lễ tân chọn khách có ưu tiên cao hơn (VIP, đặt trước) thay vì theo FIFO.

\vspace{0.5em}
\textbf{A2. Cập nhật thông tin (từ bước 3)} \\
3.1 Lễ tân cập nhật số điện thoại, số lượng khách hoặc hủy yêu cầu.

\end{minipage} \\ \hline

\textbf{Exception Flows} &
\begin{minipage}[t]{\linewidth}

\textbf{E1. Khách không phản hồi (từ bước 6)} \\
6.1 Nếu khách không phản hồi trong thời gian quy định, hệ thống chuyển trạng thái sang \texttt{Skipped}.

\vspace{0.5em}
\textbf{E2. Thời gian chờ thay đổi (từ bước 2)} \\
2.1 Nếu thời gian chờ tăng đột biến, hệ thống yêu cầu lễ tân xác nhận lại trước khi hiển thị.

\end{minipage} \\ \hline

\textbf{Business Rules} &
\begin{itemize}[leftmargin=*, noitemsep]
  \item Thứ tự mặc định là FIFO trong cùng nhóm sức chứa.
  \item Có thể override thứ tự theo ưu tiên.
  \item Chỉ gọi khách khi có bàn phù hợp với số lượng và yêu cầu.
  \item Thời gian giữ lượt gọi được cấu hình (timeout).
\end{itemize} \\ \hline

\end{longtable}



    \renewcommand{\arraystretch}{1.3}

\begin{longtable}{|p{4cm}|p{10cm}|}
\hline
\textbf{Use Case ID} & UC-RES-04 \\ \hline
\textbf{Use Case Name} & Cập nhật trạng thái bàn và giải phóng bàn \\ \hline
\textbf{Primary Actors} & Lễ tân, Phục vụ, Thu ngân \\ \hline

\textbf{Description} & Cập nhật trạng thái bàn, đóng phiên phục vụ và giải phóng bàn cho lượt khách tiếp theo. \\ \hline

\textbf{Trigger} & Khách rời bàn, hoàn tất phục vụ hoặc cần thay đổi trạng thái bàn. \\ \hline

\textbf{Preconditions} &
\begin{itemize}[leftmargin=*, noitemsep]
  \item Bàn tồn tại trong hệ thống.
  \item Người thao tác có quyền cập nhật trạng thái.
  \item Nếu có \texttt{TableSession}, hệ thống truy xuất được trạng thái hóa đơn.
\end{itemize} \\ \hline

\textbf{Postconditions} &
\begin{itemize}[leftmargin=*, noitemsep]
  \item Trạng thái bàn được cập nhật nhất quán.
  \item \texttt{TableSession} được đóng nếu đủ điều kiện.
  \item Lịch sử thay đổi được ghi nhận.
\end{itemize} \\ \hline

\textbf{Normal Flow} &
\begin{enumerate}[leftmargin=*, noitemsep]
  \item Người dùng chọn bàn trên sơ đồ bàn.
  \item Hệ thống hiển thị trạng thái bàn, phiên phục vụ và hóa đơn liên quan.
  \item Người dùng chọn trạng thái đích (\texttt{Cleaning}, \texttt{Available}, \texttt{OutOfService}).
  \item Nếu chuyển sang trạng thái giải phóng, hệ thống kiểm tra hóa đơn hoặc split còn mở.
  \item Nếu tất cả hóa đơn đã hoàn tất, hệ thống đóng \texttt{TableSession} và chuyển bàn sang \texttt{Cleaning}.
  \item Phục vụ xác nhận đã dọn bàn xong.
  \item Hệ thống chuyển bàn sang trạng thái \texttt{Available}.
  \item Hệ thống cập nhật danh sách chờ và gợi ý khách tiếp theo nếu có.
\end{enumerate} \\ \hline

\textbf{Alternative Flows} &
\begin{minipage}[t]{\linewidth}

\textbf{A1. Chuyển sang OutOfService (từ bước 3)} \\
3.1 Người dùng chuyển bàn sang trạng thái \texttt{OutOfService} để bảo trì hoặc khóa khu vực.

\vspace{0.5em}
\textbf{A2. Chuyển bàn trong khi phục vụ (từ bước 3)} \\
3.2 Nếu khách đổi bàn, hệ thống cập nhật bàn mới nhưng giữ nguyên \texttt{TableSession}.

\end{minipage} \\ \hline

\textbf{Exception Flows} &
\begin{minipage}[t]{\linewidth}

\textbf{E1. Chưa thanh toán (từ bước 4)} \\
4.1 Hệ thống từ chối giải phóng bàn và yêu cầu hoàn tất thanh toán.

\vspace{0.5em}
\textbf{E2. Xung đột cập nhật (từ bước 2)} \\
2.1 Hệ thống phát hiện cập nhật đồng thời và yêu cầu tải lại dữ liệu.

\vspace{0.5em}
\textbf{E3. Không đủ quyền (từ bước 3)} \\
3.3 Hệ thống chặn thao tác nếu người dùng không có quyền.

\end{minipage} \\ \hline

\textbf{Business Rules} &
\begin{itemize}[leftmargin=*, noitemsep]
  \item Chỉ chuyển sang \texttt{Available} khi không còn \texttt{TableSession} và hóa đơn mở.
  \item Mọi thay đổi trạng thái phải được ghi log (thời gian, người thực hiện).
  \item Không được giải phóng bàn khi còn công nợ.
\end{itemize} \\ \hline

\end{longtable}



    \renewcommand{\arraystretch}{1.3}

\begin{longtable}{|p{4cm}|p{10cm}|}
\hline
\textbf{Use Case ID} & UC-RES-05 \\ \hline
\textbf{Use Case Name} & Gửi thông báo cho khách hàng \\ \hline
\textbf{Primary Actors} & Hệ thống, Lễ tân \\ \hline

\textbf{Description} & Gửi thông báo cho khách hàng khi có sự kiện như đặt bàn, thay đổi hoặc đến lượt chờ. \\ \hline

\textbf{Trigger} & Sự kiện nghiệp vụ phát sinh hoặc lễ tân kích hoạt gửi thủ công. \\ \hline

\textbf{Preconditions} &
\begin{itemize}[leftmargin=*, noitemsep]
  \item Có thông tin liên hệ hợp lệ.
  \item Mẫu thông báo và kênh gửi đã được cấu hình.
  \item Hệ thống kết nối được dịch vụ gửi thông báo.
\end{itemize} \\ \hline

\textbf{Postconditions} &
\begin{itemize}[leftmargin=*, noitemsep]
  \item Tạo bản ghi \texttt{NotificationMessage}.
  \item Thông báo được gửi hoặc đưa vào hàng chờ.
  \item Lịch sử gửi được lưu lại.
\end{itemize} \\ \hline

\textbf{Normal Flow} &
\begin{enumerate}[leftmargin=*, noitemsep]
  \item Hệ thống nhận sự kiện cần gửi thông báo.
  \item Hệ thống xác định loại thông báo và chọn mẫu nội dung phù hợp.
  \item Hệ thống kiểm tra thông tin liên hệ và kênh gửi ưu tiên.
  \item Hệ thống cá nhân hóa nội dung thông báo.
  \item Hệ thống tạo \texttt{NotificationMessage} và gửi qua kênh đã chọn.
  \item Hệ thống ghi nhận kết quả và cập nhật trạng thái gửi.
\end{enumerate} \\ \hline

\textbf{Alternative Flows} &
\begin{minipage}[t]{\linewidth}

\textbf{A1. Gửi thủ công (từ bước 1)} \\
1.1 Lễ tân kích hoạt gửi lại thông báo cho khách.

\vspace{0.5em}
\textbf{A2. Chuyển kênh gửi (từ bước 3)} \\
3.1 Nếu kênh chính không khả dụng, hệ thống chuyển sang kênh dự phòng.

\end{minipage} \\ \hline

\textbf{Exception Flows} &
\begin{minipage}[t]{\linewidth}

\textbf{E1. Thiếu thông tin liên hệ (từ bước 3)} \\
3.2 Hệ thống đánh dấu cần xử lý thủ công.

\vspace{0.5em}
\textbf{E2. Lỗi dịch vụ gửi (từ bước 5)} \\
5.1 Hệ thống đưa thông báo vào hàng chờ retry.

\vspace{0.5em}
\textbf{E3. Bị chặn spam (từ bước 5)} \\
5.2 Hệ thống bỏ qua gửi trùng và ghi nhận lý do.

\end{minipage} \\ \hline

\textbf{Business Rules} &
\begin{itemize}[leftmargin=*, noitemsep]
  \item Mỗi thông báo phải liên kết với \texttt{Reservation} hoặc \texttt{WaitlistEntry}.
  \item Thông báo có thời hạn phải kèm thời gian phản hồi rõ ràng.
  \item Không gửi trùng thông báo trong khoảng thời gian cấu hình.
\end{itemize} \\ \hline

\end{longtable}



\renewcommand{\arraystretch}{1.3}

\begin{longtable}{|p{4cm}|p{10cm}|}
\hline
\textbf{Use Case ID} & UC-ORD-01 \\ \hline
\textbf{Use Case Name} & Tạo đơn gọi món tại bàn \\ \hline
\textbf{Primary Actors} & Phục vụ \\ \hline

\textbf{Description} & Tạo đơn gọi món cho một bàn đang phục vụ, bao gồm món, số lượng và ghi chú. \\ \hline

\textbf{Trigger} & Phục vụ chọn chức năng tạo đơn gọi món trên bàn đang hoạt động. \\ \hline

\textbf{Preconditions} &
\begin{itemize}[leftmargin=*, noitemsep]
  \item \texttt{TableSession} đang hoạt động.
  \item Phục vụ có quyền thao tác trên khu vực.
\end{itemize} \\ \hline

\textbf{Postconditions} &
\begin{itemize}[leftmargin=*, noitemsep]
  \item Đơn gọi món được lưu với trạng thái \texttt{Draft} hoặc \texttt{Confirmed}.
  \item Các dòng món được liên kết với phiên phục vụ.
\end{itemize} \\ \hline

\textbf{Normal Flow} &
\begin{enumerate}[leftmargin=*, noitemsep]
  \item Phục vụ mở bàn và chọn chức năng tạo đơn gọi món.
  \item Hệ thống kiểm tra trạng thái \texttt{TableSession}.
  \item Hệ thống hiển thị thực đơn và danh sách món.
  \item Phục vụ chọn món, số lượng và ghi chú cho từng dòng món.
  \item Hệ thống tính tạm tính theo giá hiện hành.
  \item Phục vụ lưu đơn ở trạng thái \texttt{Draft} hoặc xác nhận gửi bếp.
  \item Hệ thống lưu đơn gọi món và cập nhật trạng thái.
\end{enumerate} \\ \hline

\textbf{Alternative Flows} &
\begin{minipage}[t]{\linewidth}

\textbf{A1. Tạo nhiều đơn (từ bước 1)} \\
1.1 Phục vụ tạo nhiều đơn gọi món liên tiếp trong cùng một phiên phục vụ.

\vspace{0.5em}
\textbf{A2. Dùng mẫu đơn (từ bước 3)} \\
3.1 Phục vụ chọn mẫu đơn có sẵn cho bàn đoàn hoặc combo.

\end{minipage} \\ \hline

\textbf{Exception Flows} &
\begin{minipage}[t]{\linewidth}

\textbf{E1. Món hết hàng (từ bước 4)} \\
4.1 Hệ thống từ chối thêm món và yêu cầu chọn món thay thế.

\vspace{0.5em}
\textbf{E2. Phiên bàn không hợp lệ (từ bước 2)} \\
2.1 Hệ thống từ chối tạo đơn nếu \texttt{TableSession} đã đóng.

\end{minipage} \\ \hline

\textbf{Business Rules} &
\begin{itemize}[leftmargin=*, noitemsep]
  \item Mỗi đơn phải thuộc một \texttt{TableSession}.
  \item Giá món được cố định tại thời điểm xác nhận đơn.
  \item Một bàn có thể có nhiều đơn trong cùng phiên phục vụ.
\end{itemize} \\ \hline

\end{longtable}


   \renewcommand{\arraystretch}{1.3}

\begin{longtable}{|p{4cm}|p{10cm}|}
\hline
\textbf{Use Case ID} & UC-ORD-02 \\ \hline
\textbf{Use Case Name} & Tùy biến món và ghi chú dị ứng \\ \hline
\textbf{Primary Actors} & Phục vụ, Bếp \\ \hline

\textbf{Description} & Ghi nhận tùy chọn món, mức chế biến và ghi chú dị ứng cho từng dòng món. \\ \hline

\textbf{Trigger} & Phục vụ chọn chỉnh sửa chi tiết một dòng món trong đơn gọi món. \\ \hline

\textbf{Preconditions} &
\begin{itemize}[leftmargin=*, noitemsep]
  \item Món hỗ trợ tùy chọn hoặc ghi chú.
  \item Đơn gọi món ở trạng thái \texttt{Draft} hoặc chưa gửi bếp.
  \item Cấu hình modifier đã được thiết lập.
\end{itemize} \\ \hline

\textbf{Postconditions} &
\begin{itemize}[leftmargin=*, noitemsep]
  \item Tùy chọn được lưu dưới dạng \texttt{OrderItemModifier}.
  \item Ghi chú dị ứng được gắn vào dòng món.
\end{itemize} \\ \hline

\textbf{Normal Flow} &
\begin{enumerate}[leftmargin=*, noitemsep]
  \item Phục vụ mở chi tiết dòng món.
  \item Hệ thống hiển thị các nhóm tùy chọn và giới hạn chọn.
  \item Phục vụ chọn các tùy biến cho món.
  \item Phục vụ nhập ghi chú dị ứng hoặc yêu cầu đặc biệt.
  \item Hệ thống kiểm tra quy tắc chọn (min/max).
  \item Hệ thống lưu cấu hình tùy chọn vào dòng món.
  \item Khi đơn được gửi bếp, thông tin tùy biến và ghi chú được hiển thị.
\end{enumerate} \\ \hline

\textbf{Alternative Flows} &
\begin{minipage}[t]{\linewidth}

\textbf{A1. Dùng cấu hình mẫu (từ bước 2)} \\
2.1 Phục vụ chọn nhanh cấu hình có sẵn cho món.

\vspace{0.5em}
\textbf{A2. Gắn cờ dị ứng (từ bước 4)} \\
4.1 Hệ thống tự động đánh dấu nổi bật các ghi chú dị ứng quan trọng.

\end{minipage} \\ \hline

\textbf{Exception Flows} &
\begin{minipage}[t]{\linewidth}

\textbf{E1. Vi phạm quy tắc chọn (từ bước 5)} \\
5.1 Hệ thống yêu cầu điều chỉnh số lượng lựa chọn.

\vspace{0.5em}
\textbf{E2. Ghi chú quá dài (từ bước 4)} \\
4.2 Hệ thống yêu cầu rút gọn hoặc chuyển sang ghi chú nội bộ.

\vspace{0.5em}
\textbf{E3. Đơn đã gửi bếp (từ bước 1)} \\
1.1 Hệ thống từ chối chỉnh sửa nếu đơn đã được xác nhận.

\end{minipage} \\ \hline

\textbf{Business Rules} &
\begin{itemize}[leftmargin=*, noitemsep]
  \item Ghi chú dị ứng phải hiển thị rõ ràng trên phiếu bếp.
  \item Modifier có thể thay đổi giá món.
  \item Không cho phép chỉnh sửa sau khi gửi bếp.
\end{itemize} \\ \hline

\end{longtable}



   \renewcommand{\arraystretch}{1.3}

\begin{longtable}{|p{4cm}|p{10cm}|}
\hline
\textbf{Use Case ID} & UC-ORD-03 \\ \hline
\textbf{Use Case Name} & Gửi đơn gọi món sang bếp \\ \hline
\textbf{Primary Actors} & Phục vụ, Bếp \\ \hline

\textbf{Description} & Chuyển đơn gọi món đã xác nhận đến các trạm bếp để chế biến. \\ \hline

\textbf{Trigger} & Phục vụ nhấn gửi bếp hoặc hệ thống tự động gửi. \\ \hline

\textbf{Preconditions} &
\begin{itemize}[leftmargin=*, noitemsep]
  \item Đơn có ít nhất một dòng món hợp lệ.
  \item Đơn ở trạng thái \texttt{Draft}.
  \item Kết nối tới hệ thống bếp khả dụng.
\end{itemize} \\ \hline

\textbf{Postconditions} &
\begin{itemize}[leftmargin=*, noitemsep]
  \item Tạo một hoặc nhiều \texttt{KitchenTicket}.
  \item Đơn chuyển sang trạng thái \texttt{Confirmed/Queued}.
\end{itemize} \\ \hline

\textbf{Normal Flow} &
\begin{enumerate}[leftmargin=*, noitemsep]
  \item Phục vụ xác nhận các dòng món cần gửi.
  \item Hệ thống kiểm tra trạng thái món, modifier và tồn kho logic.
  \item Hệ thống xác định trạm bếp cho từng dòng món.
  \item Hệ thống tách đơn thành các \texttt{KitchenTicket} theo trạm bếp.
  \item Hệ thống xếp hàng phiếu bếp theo ưu tiên và thời gian.
  \item Hệ thống cập nhật trạng thái đơn thành \texttt{Confirmed/Queued}.
  \item Màn hình bếp hiển thị phiếu và phát tín hiệu thông báo.
\end{enumerate} \\ \hline

\textbf{Alternative Flows} &
\begin{minipage}[t]{\linewidth}

\textbf{A1. Gửi món sau (từ bước 1)} \\
1.1 Phục vụ đánh dấu một số món để gửi sau.

\vspace{0.5em}
\textbf{A2. Gộp phiếu bếp (từ bước 4)} \\
4.1 Hệ thống gộp các phiếu bếp gần nhau theo cấu hình thời gian.

\end{minipage} \\ \hline

\textbf{Exception Flows} &
\begin{minipage}[t]{\linewidth}

\textbf{E1. Không có trạm bếp (từ bước 3)} \\
3.1 Hệ thống chặn gửi và cảnh báo lỗi cấu hình menu.

\vspace{0.5em}
\textbf{E2. Mất kết nối bếp (từ bước 5)} \\
5.1 Hệ thống đưa phiếu vào hàng đợi và retry sau.

\end{minipage} \\ \hline

\textbf{Business Rules} &
\begin{itemize}[leftmargin=*, noitemsep]
  \item Mỗi dòng món phải được định tuyến đến ít nhất một trạm bếp.
  \item Chỉ dòng món hợp lệ mới được gửi.
  \item Không cho phép gửi lại các dòng đã được gửi trước đó.
\end{itemize} \\ \hline

\end{longtable}



    \renewcommand{\arraystretch}{1.3}

\begin{longtable}{|p{4cm}|p{10cm}|}
\hline
\textbf{Use Case ID} & UC-KIT-01 \\ \hline
\textbf{Use Case Name} & Cập nhật tiến độ chế biến \\ \hline
\textbf{Primary Actors} & Bếp \\ \hline

\textbf{Description} & Theo dõi và cập nhật trạng thái món từ khi nhận phiếu bếp đến khi sẵn sàng phục vụ. \\ \hline

\textbf{Trigger} & Nhân viên bếp chọn một phiếu bếp trên màn hình KDS. \\ \hline

\textbf{Preconditions} &
\begin{itemize}[leftmargin=*, noitemsep]
  \item \texttt{KitchenTicket} đã được tạo và gán trạm bếp.
  \item Nhân viên đã đăng nhập vào hệ thống bếp.
\end{itemize} \\ \hline

\textbf{Postconditions} &
\begin{itemize}[leftmargin=*, noitemsep]
  \item Trạng thái món và phiếu bếp được cập nhật.
  \item POS nhận được tiến độ mới nhất.
\end{itemize} \\ \hline

\textbf{Normal Flow} &
\begin{enumerate}[leftmargin=*, noitemsep]
  \item Nhân viên bếp mở phiếu bếp ở trạng thái \texttt{Queued}.
  \item Nhân viên cập nhật trạng thái từng món (\texttt{Started}, \texttt{Cooking}, \texttt{Ready}).
  \item Hệ thống ghi nhận thời gian cho từng trạng thái.
  \item Hệ thống so sánh thời gian chế biến với SLA.
  \item Khi tất cả món hoàn tất, hệ thống chuyển phiếu sang \texttt{Ready}.
  \item Hệ thống gửi thông báo đến POS để phục vụ ra món.
\end{enumerate} \\ \hline

\textbf{Alternative Flows} &
\begin{minipage}[t]{\linewidth}

\textbf{A1. Ưu tiên phiếu (từ bước 1)} \\
1.1 Nhân viên bếp ưu tiên xử lý phiếu VIP hoặc món cần ra trước.

\vspace{0.5em}
\textbf{A2. Giữ món (từ bước 2)} \\
2.1 Nhân viên đánh dấu \texttt{Hold for service} để chờ đồng bộ với món khác.

\end{minipage} \\ \hline

\textbf{Exception Flows} &
\begin{minipage}[t]{\linewidth}

\textbf{E1. Thiếu nguyên liệu (từ bước 2)} \\
2.2 Nhân viên đánh dấu \texttt{Blocked} và gửi yêu cầu thay món.

\vspace{0.5em}
\textbf{E2. Mất kết nối KDS (từ bước 3)} \\
3.1 Hệ thống chuyển sang hàng đợi cục bộ và đồng bộ lại sau.

\end{minipage} \\ \hline

\textbf{Business Rules} &
\begin{itemize}[leftmargin=*, noitemsep]
  \item Chỉ trạm bếp được phân công mới được cập nhật trạng thái.
  \item Mọi thay đổi trạng thái phải được ghi log.
  \item Trạng thái phiếu phụ thuộc vào trạng thái tất cả món.
\end{itemize} \\ \hline

\end{longtable}



    \renewcommand{\arraystretch}{1.3}

\begin{longtable}{|p{4cm}|p{10cm}|}
\hline
\textbf{Use Case ID} & UC-KIT-02 \\ \hline
\textbf{Use Case Name} & Ra món cho bàn \\ \hline
\textbf{Primary Actors} & Phục vụ \\ \hline

\textbf{Description} & Nhận món từ bếp và phục vụ đúng bàn, đúng thứ tự. \\ \hline

\textbf{Trigger} & Hệ thống thông báo món hoặc phiếu bếp ở trạng thái \texttt{Ready}. \\ \hline

\textbf{Preconditions} &
\begin{itemize}[leftmargin=*, noitemsep]
  \item Món ở trạng thái \texttt{Ready}.
  \item \texttt{TableSession} còn hoạt động.
\end{itemize} \\ \hline

\textbf{Postconditions} &
\begin{itemize}[leftmargin=*, noitemsep]
  \item Dòng món được cập nhật trạng thái \texttt{Served}.
  \item Thời gian phục vụ được ghi nhận.
\end{itemize} \\ \hline

\textbf{Normal Flow} &
\begin{enumerate}[leftmargin=*, noitemsep]
  \item Phục vụ mở danh sách món sẵn sàng.
  \item Phục vụ xác nhận nhận món từ bếp.
  \item Hệ thống hiển thị thông tin bàn, vị trí ngồi và ghi chú.
  \item Phục vụ giao món cho khách.
  \item Phục vụ xác nhận \enquote{Đã phục vụ}.
  \item Hệ thống cập nhật trạng thái từng dòng món thành \texttt{Served}.
  \item Nếu tất cả món đã phục vụ, hệ thống cập nhật trạng thái phục vụ của bàn.
\end{enumerate} \\ \hline

\textbf{Alternative Flows} &
\begin{minipage}[t]{\linewidth}

\textbf{A1. Xác nhận theo phiếu (từ bước 1)} \\
1.1 Phục vụ xác nhận toàn bộ phiếu thay vì từng món.

\vspace{0.5em}
\textbf{A2. Giữ món theo course (từ bước 3)} \\
3.1 Phục vụ trì hoãn ra món để đồng bộ theo course.

\end{minipage} \\ \hline

\textbf{Exception Flows} &
\begin{minipage}[t]{\linewidth}

\textbf{E1. Sai món hoặc sai tùy biến (từ bước 3)} \\
3.2 Phục vụ trả món về bếp và ghi lý do.

\vspace{0.5em}
\textbf{E2. Bàn tạm vắng (từ bước 4)} \\
4.1 Phục vụ đánh dấu trì hoãn phục vụ.

\vspace{0.5em}
\textbf{E3. Trạng thái không hợp lệ (từ bước 5)} \\
5.1 Hệ thống từ chối nếu món chưa ở trạng thái \texttt{Ready}.

\end{minipage} \\ \hline

\textbf{Business Rules} &
\begin{itemize}[leftmargin=*, noitemsep]
  \item Chỉ món ở trạng thái \texttt{Ready} mới được phục vụ.
  \item Mọi thao tác trả món phải có lý do.
  \item Trạng thái phục vụ phải được ghi log để phân tích SLA.
\end{itemize} \\ \hline

\end{longtable}



    \renewcommand{\arraystretch}{1.3}

\begin{longtable}{|p{4cm}|p{10cm}|}
\hline
\textbf{Use Case ID} & UC-KIT-03 \\ \hline
\textbf{Use Case Name} & Ưu tiên hàng chờ bếp và xử lý món gấp \\ \hline
\textbf{Primary Actors} & Bếp, Quản lý \\ \hline

\textbf{Description} & Điều chỉnh mức ưu tiên của ticket hoặc món để xử lý các trường hợp gấp và giảm trễ hạn. \\ \hline

\textbf{Trigger} & Ticket gần vượt SLA, khách VIP hoặc yêu cầu can thiệp từ quản lý. \\ \hline

\textbf{Preconditions} &
\begin{itemize}[leftmargin=*, noitemsep]
  \item \texttt{KitchenTicket} hoặc item tồn tại trong hàng chờ.
  \item Có dữ liệu SLA và thứ tự hàng chờ.
  \item Người thao tác có quyền thay đổi ưu tiên.
\end{itemize} \\ \hline

\textbf{Postconditions} &
\begin{itemize}[leftmargin=*, noitemsep]
  \item Mức ưu tiên được cập nhật.
  \item Hàng chờ được sắp xếp lại.
  \item Lý do thay đổi được ghi log.
\end{itemize} \\ \hline

\textbf{Normal Flow} &
\begin{enumerate}[leftmargin=*, noitemsep]
  \item Nhân viên mở hàng chờ của trạm bếp.
  \item Hệ thống hiển thị độ trễ và các ticket có nguy cơ trễ.
  \item Người thao tác chọn ticket hoặc món cần ưu tiên.
  \item Hệ thống áp dụng quy tắc ưu tiên (SLA, VIP, course).
  \item Người thao tác xác nhận hoặc nhập lý do ưu tiên.
  \item Hệ thống cập nhật mức ưu tiên.
  \item Hệ thống sắp xếp lại hàng chờ.
  \item Hệ thống làm nổi bật ticket ưu tiên trên KDS.
\end{enumerate} \\ \hline

\textbf{Alternative Flows} &
\begin{minipage}[t]{\linewidth}

\textbf{A1. Tự động ưu tiên (từ bước 2)} \\
2.1 Hệ thống tự động tăng ưu tiên nếu ticket gần vượt SLA.

\vspace{0.5em}
\textbf{A2. Ưu tiên theo món (từ bước 3)} \\
3.1 Người thao tác chỉ ưu tiên một phần của ticket (ví dụ món khai vị).

\end{minipage} \\ \hline

\textbf{Exception Flows} &
\begin{minipage}[t]{\linewidth}

\textbf{E1. Không đủ quyền (từ bước 3)} \\
3.2 Hệ thống yêu cầu xác nhận từ quản lý.

\vspace{0.5em}
\textbf{E2. Quá tải trạm bếp (từ bước 2)} \\
2.2 Hệ thống cảnh báo và gợi ý điều phối sang trạm khác.

\vspace{0.5em}
\textbf{E3. Thiếu dữ liệu SLA (từ bước 4)} \\
4.1 Hệ thống chỉ cho phép ưu tiên thủ công.

\end{minipage} \\ \hline

\textbf{Business Rules} &
\begin{itemize}[leftmargin=*, noitemsep]
  \item Mọi hành động ưu tiên phải có lý do.
  \item Không được gây trì hoãn vô hạn cho ticket khác (anti-starvation).
  \item Ưu tiên có thể áp dụng ở mức ticket hoặc item.
\end{itemize} \\ \hline

\end{longtable}


    \renewcommand{\arraystretch}{1.3}

\begin{longtable}{|p{4cm}|p{10cm}|}
\hline
\textbf{Use Case ID} & UC-ORD-04 \\ \hline
\textbf{Use Case Name} & Sửa hoặc hủy dòng món \\ \hline
\textbf{Primary Actors} & Phục vụ, Quản lý \\ \hline

\textbf{Description} & Điều chỉnh hoặc hủy dòng món trong đơn gọi món theo trạng thái và quyền hạn. \\ \hline

\textbf{Trigger} & Phục vụ chọn dòng món để sửa hoặc hủy. \\ \hline

\textbf{Preconditions} &
\begin{itemize}[leftmargin=*, noitemsep]
  \item Đơn chưa hoàn tất thanh toán.
  \item Người dùng có quyền chỉnh sửa theo trạng thái món.
\end{itemize} \\ \hline

\textbf{Postconditions} &
\begin{itemize}[leftmargin=*, noitemsep]
  \item Dòng món được cập nhật hoặc chuyển trạng thái \texttt{Voided}.
  \item Lịch sử thay đổi được ghi nhận.
  \item Nếu đã gửi bếp, hệ thống phát thông báo đến KDS.
\end{itemize} \\ \hline

\textbf{Normal Flow} &
\begin{enumerate}[leftmargin=*, noitemsep]
  \item Phục vụ chọn dòng món cần chỉnh sửa.
  \item Hệ thống hiển thị trạng thái món và quyền chỉnh sửa.
  \item Phục vụ chọn sửa thông tin hoặc hủy dòng món.
  \item Nếu món đã gửi bếp, hệ thống yêu cầu nhập lý do.
  \item Nếu cần, hệ thống yêu cầu phê duyệt từ quản lý.
  \item Hệ thống cập nhật dữ liệu dòng món.
  \item Hệ thống phát sự kiện đến KDS và cập nhật tạm tính.
\end{enumerate} \\ \hline

\textbf{Alternative Flows} &
\begin{minipage}[t]{\linewidth}

\textbf{A1. Chỉnh sửa trước khi gửi bếp (từ bước 3)} \\
3.1 Phục vụ chỉnh sửa tự do khi món ở trạng thái \texttt{Draft}.

\vspace{0.5em}
\textbf{A2. Hủy sau khi chế biến (từ bước 4)} \\
4.1 Hệ thống yêu cầu phê duyệt bắt buộc khi món đã ở trạng thái \texttt{Cooking} hoặc \texttt{Ready}.

\end{minipage} \\ \hline

\textbf{Exception Flows} &
\begin{minipage}[t]{\linewidth}

\textbf{E1. Không đủ quyền (từ bước 5)} \\
5.1 Hệ thống từ chối thao tác nếu không có quyền.

\vspace{0.5em}
\textbf{E2. Xung đột cập nhật (từ bước 2)} \\
2.1 Hệ thống phát hiện chỉnh sửa đồng thời và yêu cầu tải lại.

\end{minipage} \\ \hline

\textbf{Business Rules} &
\begin{itemize}[leftmargin=*, noitemsep]
  \item Hủy món sau khi gửi bếp phải có lý do.
  \item Có thể yêu cầu phê duyệt nếu vượt ngưỡng cấu hình.
  \item Không được mất lịch sử trạng thái và giá ban đầu.
\end{itemize} \\ \hline

\end{longtable}



    \begin{usecasebox}[title={UC-BIL-01 --- Tạo hóa đơn cho bàn}]
    \textbf{Tác nhân chính:} Phục vụ, Thu ngân\\
    \textbf{Mục tiêu:} Tổng hợp các đơn gọi món của bàn thành hóa đơn có thể thanh toán.\\
    \textbf{Kích hoạt:} Khách yêu cầu tính tiền hoặc phục vụ chọn \enquote{Tạo hóa đơn}.

    \medskip
    \textbf{Tiền điều kiện}
    \begin{itemize}[leftmargin=1.8em]
  \item Phiên bàn đang mở và có ít nhất một đơn gọi món hợp lệ.
  \item Các món cần tính tiền đã ở trạng thái cho phép xuất hóa đơn.
\end{itemize}

    \textbf{Hậu điều kiện}
    \begin{itemize}[leftmargin=1.8em]
  \item Một hóa đơn được tạo với tổng tiền, phí, giảm giá và trạng thái thanh toán ban đầu.
  \item Bàn có thể được chuyển sang trạng thái chờ thanh toán.
\end{itemize}

    \textbf{Luồng chính}
    \begin{enumerate}[leftmargin=1.8em]
  \item Phục vụ chọn bàn cần tính tiền.
  \item Hệ thống gom toàn bộ đơn gọi món chưa được quyết toán của bàn.
  \item Hệ thống tính subtotal, phụ phí, thuế, khuyến mãi và tổng thanh toán.
  \item Phục vụ/thu ngân xem trước hóa đơn và xác nhận phát hành.
  \item Hệ thống tạo bản ghi \texttt{Bill} và danh sách phân bổ thanh toán mặc định.
\end{enumerate}

    \textbf{Luồng thay thế}
    \begin{itemize}[leftmargin=1.8em]
  \item Nhà hàng có thể tạo hóa đơn tạm thời để khách kiểm tra trước khi thanh toán.
  \item Có thể tạo nhiều hóa đơn con nếu bàn yêu cầu tách theo ghế hoặc theo nhóm món.
\end{itemize}

    \textbf{Ngoại lệ}
    \begin{itemize}[leftmargin=1.8em]
  \item Một đơn gọi món vừa bị hủy hoặc điều chỉnh trong lúc tạo hóa đơn: hệ thống yêu cầu tính lại.
  \item Chưa có cấu hình thuế hoặc phí cho chi nhánh: hệ thống chặn phát hành hóa đơn.
\end{itemize}

    \textbf{Quy tắc nghiệp vụ}
    \begin{itemize}[leftmargin=1.8em]
  \item Một đơn gọi món không được gắn vào hơn một hóa đơn đang mở cùng lúc.
  \item Công thức tính phí và thuế phải được chụp lại tại thời điểm tạo hóa đơn.
\end{itemize}
    \end{usecasebox}



    \begin{usecasebox}[title={UC-BIL-02 --- Tách hóa đơn và thêm tiền boa}]
    \textbf{Tác nhân chính:} Thu ngân, Phục vụ\\
    \textbf{Mục tiêu:} Chia hóa đơn thành nhiều phần thanh toán và ghi nhận tiền boa.\\
    \textbf{Kích hoạt:} Khách yêu cầu thanh toán riêng theo đầu người, theo món hoặc theo số tiền.

    \medskip
    \textbf{Tiền điều kiện}
    \begin{itemize}[leftmargin=1.8em]
  \item Hóa đơn đã được tạo.
  \item Hóa đơn chưa đóng hoàn toàn.
\end{itemize}

    \textbf{Hậu điều kiện}
    \begin{itemize}[leftmargin=1.8em]
  \item Các phần thanh toán (\texttt{BillSplit}) được tạo và sẵn sàng nhận giao dịch thanh toán.
  \item Tiền boa được phân bổ theo chính sách nhà hàng.
\end{itemize}

    \textbf{Luồng chính}
    \begin{enumerate}[leftmargin=1.8em]
  \item Thu ngân mở hóa đơn cần tách.
  \item Hệ thống hiển thị các món, số ghế, subtotal và các quy tắc chia hóa đơn.
  \item Thu ngân chọn kiểu chia: đều, theo món, theo ghế, theo số tiền hoặc kết hợp.
  \item Thu ngân nhập tiền boa hoặc chọn mức tiền boa gợi ý.
  \item Hệ thống tính lại số tiền của từng phần thanh toán và hiển thị chênh lệch nếu có.
  \item Thu ngân xác nhận cấu hình tách hóa đơn.
\end{enumerate}

    \textbf{Luồng thay thế}
    \begin{itemize}[leftmargin=1.8em]
  \item Khách có thể chỉ tách một phần món ra thanh toán trước, phần còn lại giữ nguyên hóa đơn gốc.
  \item Tiền boa có thể áp dụng cho toàn bộ hóa đơn hoặc gắn riêng từng split.
\end{itemize}

    \textbf{Ngoại lệ}
    \begin{itemize}[leftmargin=1.8em]
  \item Tổng các split không khớp tổng hóa đơn: hệ thống từ chối lưu và yêu cầu cân bằng lại.
  \item Một split đã thanh toán không được phép thay đổi phạm vi món trừ khi hoàn tiền/phê duyệt.
\end{itemize}

    \textbf{Quy tắc nghiệp vụ}
    \begin{itemize}[leftmargin=1.8em]
  \item Tổng các split cộng tiền boa phải khớp với công thức tiền của hóa đơn.
  \item Tiền boa có thể bị giới hạn tối đa theo chính sách địa phương hoặc hệ thống POS.
\end{itemize}
    \end{usecasebox}



    \begin{usecasebox}[title={UC-PAY-01 --- Xử lý thanh toán}]
    \textbf{Tác nhân chính:} Thu ngân\\
    \textbf{Mục tiêu:} Thu tiền cho một hóa đơn hoặc một split bằng tiền mặt, thẻ hoặc ví điện tử.\\
    \textbf{Kích hoạt:} Thu ngân chọn \enquote{Thanh toán} trên hóa đơn hoặc split cụ thể.

    \medskip
    \textbf{Tiền điều kiện}
    \begin{itemize}[leftmargin=1.8em]
  \item Hóa đơn/split ở trạng thái chờ thanh toán.
  \item Cổng thanh toán và cấu hình máy POS khả dụng đối với phương thức không tiền mặt.
\end{itemize}

    \textbf{Hậu điều kiện}
    \begin{itemize}[leftmargin=1.8em]
  \item Một bản ghi \texttt{Payment} được tạo với trạng thái thành công, thất bại hoặc chờ xác nhận.
  \item Nếu đã thanh toán đủ, hóa đơn được đóng.
\end{itemize}

    \textbf{Luồng chính}
    \begin{enumerate}[leftmargin=1.8em]
  \item Thu ngân chọn split cần thanh toán và phương thức thanh toán.
  \item Hệ thống hiển thị số tiền phải thu và các bước xác nhận.
  \item Nếu thanh toán điện tử, hệ thống gọi \texttt{PaymentGateway} để ủy quyền và ghi nhận giao dịch.
  \item Nếu thanh toán tiền mặt, thu ngân nhập số tiền khách đưa và hệ thống tính tiền thối.
  \item Hệ thống lưu kết quả thanh toán, cập nhật số dư còn lại và trạng thái hóa đơn.
  \item Nếu hóa đơn đã được thanh toán đủ, hệ thống đóng hóa đơn và mở bước phát hành biên lai.
\end{enumerate}

    \textbf{Luồng thay thế}
    \begin{itemize}[leftmargin=1.8em]
  \item Một hóa đơn có thể được thanh toán bằng nhiều phương thức trong nhiều lần khác nhau.
  \item Nhà hàng có thể cho phép đặt cọc trước và trừ vào tổng hóa đơn khi check-out.
\end{itemize}

    \textbf{Ngoại lệ}
    \begin{itemize}[leftmargin=1.8em]
  \item Cổng thanh toán trả về thất bại hoặc quá thời gian chờ: hệ thống ghi nhận trạng thái \texttt{Failed/Pending Review}.
  \item Thiết bị thanh toán ngoại vi không phản hồi: hệ thống cho phép chuyển phương thức khác.
  \item Số tiền thanh toán vượt mức còn lại: hệ thống yêu cầu xác nhận hoặc tự giới hạn theo cấu hình.
\end{itemize}

    \textbf{Quy tắc nghiệp vụ}
    \begin{itemize}[leftmargin=1.8em]
  \item Mọi thanh toán điện tử phải lưu mã tham chiếu giao dịch.
  \item Hóa đơn chỉ đóng khi tổng thanh toán hợp lệ lớn hơn hoặc bằng tổng phải thu.
\end{itemize}
    \end{usecasebox}



    \begin{usecasebox}[title={UC-PAY-02 --- Phát hành biên lai}]
    \textbf{Tác nhân chính:} Thu ngân, Khách hàng\\
    \textbf{Mục tiêu:} Tạo và gửi biên lai sau khi thanh toán thành công.\\
    \textbf{Kích hoạt:} Một hóa đơn hoặc split được thanh toán thành công.

    \medskip
    \textbf{Tiền điều kiện}
    \begin{itemize}[leftmargin=1.8em]
  \item Thanh toán đã được ghi nhận hợp lệ.
  \item Kênh phát hành biên lai (in giấy, email, SMS) đã được cấu hình.
\end{itemize}

    \textbf{Hậu điều kiện}
    \begin{itemize}[leftmargin=1.8em]
  \item Một \texttt{Receipt} được tạo và liên kết với hóa đơn/thanh toán.
  \item Khách hàng nhận biên lai theo kênh đã chọn.
\end{itemize}

    \textbf{Luồng chính}
    \begin{enumerate}[leftmargin=1.8em]
  \item Hệ thống tạo nội dung biên lai gồm món, giảm giá, phí, phương thức thanh toán và thời gian.
  \item Thu ngân chọn in giấy hoặc gửi điện tử.
  \item Hệ thống phát hành biên lai và cập nhật trạng thái thành công.
  \item Nếu cần, hệ thống gửi bản sao tới email khách hàng hoặc số điện thoại đã lưu.
\end{enumerate}

    \textbf{Luồng thay thế}
    \begin{itemize}[leftmargin=1.8em]
  \item Khách có thể yêu cầu xuất biên lai rút gọn hoặc chi tiết theo chính sách nhà hàng.
  \item Trong trường hợp split bill, mỗi split có thể có biên lai riêng.
\end{itemize}

    \textbf{Ngoại lệ}
    \begin{itemize}[leftmargin=1.8em]
  \item Máy in lỗi hoặc hết giấy: hệ thống gợi ý chuyển sang gửi điện tử hoặc in lại sau.
  \item Email/SMS không gửi được: hệ thống vẫn lưu biên lai và cho phép gửi lại thủ công.
\end{itemize}

    \textbf{Quy tắc nghiệp vụ}
    \begin{itemize}[leftmargin=1.8em]
  \item Biên lai phải tham chiếu tới ít nhất một thanh toán thành công.
  \item Việc phát hành lại biên lai không được tạo thêm thanh toán mới.
\end{itemize}
    \end{usecasebox}



    \begin{usecasebox}[title={UC-PAY-03 --- Hoàn tiền}]
    \textbf{Tác nhân chính:} Thu ngân, Quản lý\\
    \textbf{Mục tiêu:} Thực hiện hoàn tiền một phần hoặc toàn phần cho giao dịch đã thanh toán.\\
    \textbf{Kích hoạt:} Khách khiếu nại, hủy món sau thanh toán hoặc cần điều chỉnh hóa đơn.

    \medskip
    \textbf{Tiền điều kiện}
    \begin{itemize}[leftmargin=1.8em]
  \item Đã tồn tại thanh toán thành công cần hoàn.
  \item Người thao tác có quyền hoàn tiền hoặc được quản lý phê duyệt.
\end{itemize}

    \textbf{Hậu điều kiện}
    \begin{itemize}[leftmargin=1.8em]
  \item Một bản ghi \texttt{Refund} được tạo và liên kết với giao dịch thanh toán gốc.
  \item Sổ audit ghi nhận người hoàn, lý do và số tiền hoàn.
\end{itemize}

    \textbf{Luồng chính}
    \begin{enumerate}[leftmargin=1.8em]
  \item Thu ngân chọn giao dịch thanh toán cần hoàn.
  \item Hệ thống hiển thị số tiền đã thanh toán, số tiền đã hoàn trước đó và mức hoàn còn khả dụng.
  \item Thu ngân nhập số tiền hoàn và lý do.
  \item Nếu vượt ngưỡng cho phép, hệ thống yêu cầu quản lý xác nhận.
  \item Hệ thống gửi yêu cầu hoàn tiền tới gateway hoặc ghi nhận hoàn tiền mặt.
  \item Hệ thống cập nhật trạng thái hoàn tiền, số dư hóa đơn và nhật ký kiểm toán.
\end{enumerate}

    \textbf{Luồng thay thế}
    \begin{itemize}[leftmargin=1.8em]
  \item Có thể hoàn tiền một phần cho riêng món bị lỗi mà không đóng lại toàn bộ hóa đơn.
  \item Quản lý có thể chuyển hoàn tiền thành tín dụng khách hàng nếu nhà hàng áp dụng chính sách ví nội bộ.
\end{itemize}

    \textbf{Ngoại lệ}
    \begin{itemize}[leftmargin=1.8em]
  \item Gateway từ chối hoàn tiền: hệ thống giữ trạng thái \texttt{Pending Review} và cảnh báo kế toán.
  \item Số tiền hoàn vượt mức thanh toán ròng còn lại: hệ thống chặn thao tác.
  \item Giao dịch thanh toán gốc không cho phép hoàn tiền sau thời hạn quy định: hệ thống yêu cầu xử lý ngoại lệ.
\end{itemize}

    \textbf{Quy tắc nghiệp vụ}
    \begin{itemize}[leftmargin=1.8em]
  \item Mỗi khoản hoàn tiền phải tham chiếu tới giao dịch thanh toán gốc và lý do hoàn tiền chuẩn hóa.
  \item Refund vượt ngưỡng cấu hình phải có phê duyệt của quản lý.
\end{itemize}
    \end{usecasebox}



    \begin{usecasebox}[title={UC-INV-01 --- Cập nhật tiêu hao tồn kho sau phục vụ}]
    \textbf{Tác nhân chính:} Hệ thống, Quản lý kho\\
    \textbf{Mục tiêu:} Khấu trừ tồn kho theo món đã bán hoặc điều chỉnh thủ công khi cần.\\
    \textbf{Kích hoạt:} Dòng món được xác nhận đã phục vụ hoặc quản lý kho tạo giao dịch điều chỉnh.

    \medskip
    \textbf{Tiền điều kiện}
    \begin{itemize}[leftmargin=1.8em]
  \item Món ăn đã được ánh xạ công thức nguyên liệu qua \texttt{RecipeLine}.
  \item Tồn kho của chi nhánh đang hoạt động.
\end{itemize}

    \textbf{Hậu điều kiện}
    \begin{itemize}[leftmargin=1.8em]
  \item Các giao dịch kho (\texttt{StockTransaction}) được ghi nhận.
  \item Số lượng tồn (\texttt{onHand}) của nguyên liệu được cập nhật.
\end{itemize}

    \textbf{Luồng chính}
    \begin{enumerate}[leftmargin=1.8em]
  \item Khi món được phục vụ, hệ thống tra công thức nguyên liệu tương ứng.
  \item Hệ thống tính tổng định mức theo số lượng món đã ra.
  \item Hệ thống tạo các giao dịch kho trừ tồn cho từng nguyên liệu.
  \item Nếu có điều chỉnh thủ công, quản lý kho nhập số lượng và lý do điều chỉnh.
  \item Hệ thống cập nhật tồn kho hiện tại và lịch sử giao dịch.
\end{enumerate}

    \textbf{Luồng thay thế}
    \begin{itemize}[leftmargin=1.8em]
  \item Nhà hàng có thể cấu hình khấu trừ khi gửi bếp thay vì khi phục vụ để phản ánh thực tế chế biến.
  \item Các nguyên liệu hao hụt do prep đầu ca được ghi nhận bằng giao dịch riêng.
\end{itemize}

    \textbf{Ngoại lệ}
    \begin{itemize}[leftmargin=1.8em]
  \item Thiếu ánh xạ công thức: hệ thống không tự trừ tồn và tạo cảnh báo cho quản lý menu/kho.
  \item Cập nhật đồng thời dẫn đến âm kho bất thường: hệ thống chặn hoặc đưa vào hàng chờ rà soát.
\end{itemize}

    \textbf{Quy tắc nghiệp vụ}
    \begin{itemize}[leftmargin=1.8em]
  \item Mọi điều chỉnh thủ công phải có lý do chuẩn hóa như waste, spoilage, recount, prep.
  \item Tồn kho và lịch sử giao dịch phải được tách theo chi nhánh.
\end{itemize}
    \end{usecasebox}



    \begin{usecasebox}[title={UC-INV-02 --- Nhận cảnh báo sắp hết hàng}]
    \textbf{Tác nhân chính:} Quản lý, Bếp, Hệ thống\\
    \textbf{Mục tiêu:} Phát hiện nguyên liệu xuống dưới ngưỡng và gửi cảnh báo vận hành kịp thời.\\
    \textbf{Kích hoạt:} Tồn kho sau cập nhật nhỏ hơn ngưỡng trong \texttt{LowStockRule}.

    \medskip
    \textbf{Tiền điều kiện}
    \begin{itemize}[leftmargin=1.8em]
  \item Ngưỡng cảnh báo tồn kho đã được cấu hình cho từng nguyên liệu và chi nhánh.
  \item Kênh gửi cảnh báo nội bộ hoạt động.
\end{itemize}

    \textbf{Hậu điều kiện}
    \begin{itemize}[leftmargin=1.8em]
  \item Cảnh báo sắp hết hàng được gửi đến đúng vai trò chịu trách nhiệm.
  \item Người dùng có thể xem danh sách nguyên liệu cần  hoặc khóa bán món liên quan.
\end{itemize}

    \textbf{Luồng chính}
    \begin{enumerate}[leftmargin=1.8em]
  \item Sau mỗi lần cập nhật tồn, hệ thống kiểm tra điều kiện sắp hết hàng.
  \item Khi vượt ngưỡng cảnh báo, hệ thống tạo sự kiện và thông báo nội bộ.
  \item Quản lý mở danh sách cảnh báo, xem nguyên liệu ảnh hưởng và các món liên quan.
  \item Quản lý quyết định đặt hàng , điều chỉnh ngưỡng hoặc tạm ẩn món phụ thuộc.
\end{enumerate}

    \textbf{Luồng thay thế}
    \begin{itemize}[leftmargin=1.8em]
  \item Nếu nhà hàng dùng tích hợp nhà cung cấp, hệ thống có thể gợi ý lượng đặt hàng theo lịch sử tiêu thụ.
  \item Cảnh báo có thể được gom theo ca hoặc gửi ngay lập tức với nguyên liệu trọng yếu.
\end{itemize}

    \textbf{Ngoại lệ}
    \begin{itemize}[leftmargin=1.8em]
  \item Ngưỡng cấu hình thiếu hoặc sai đơn vị: hệ thống ghi cảnh báo cấu hình thay vì gửi cảnh báo vận hành.
  \item Kênh thông báo lỗi: hệ thống lưu thông báo chờ gửi lại để thực hiện lại sau.
\end{itemize}

    \textbf{Quy tắc nghiệp vụ}
    \begin{itemize}[leftmargin=1.8em]
  \item Một cảnh báo đang mở không được gửi lặp lại quá thường xuyên trong cùng khoảng thời gian chống spam.
  \item Nguyên liệu trọng yếu có thể kích hoạt tự động ẩn món phụ thuộc theo chính sách chi nhánh.
\end{itemize}
    \end{usecasebox}



    \begin{usecasebox}[title={UC-INV-03 --- Phân tích tiêu hao và đề xuất nhập hàng}]
    \textbf{Tác nhân chính:} Quản lý, Quản lý kho\\
    \textbf{Mục tiêu:} Phân tích lịch sử tiêu hao, hao hụt và cảnh báo tồn để đề xuất lượng nhập hàng phù hợp cho kỳ tiếp theo.\\
    \textbf{Kích hoạt:} Quản lý mở màn hình phân tích nhập hàng theo ngày, tuần hoặc kỳ vận hành.

    \medskip
    \textbf{Tiền điều kiện}
    \begin{itemize}[leftmargin=1.8em]
  \item Hệ thống đã lưu lịch sử \texttt{StockTransaction}, cảnh báo sắp hết hàng và số liệu tiêu hao theo chi nhánh.
  \item Quy tắc đặt lại hàng như lead time, mức an toàn hoặc min-max đã được cấu hình ở mức tối thiểu.
  \item Người dùng có quyền xem dữ liệu tồn kho và báo cáo vận hành.
\end{itemize}

    \textbf{Hậu điều kiện}
    \begin{itemize}[leftmargin=1.8em]
  \item Một tập đề xuất nhập hàng được sinh ra với lượng khuyến nghị và mức độ tin cậy.
  \item Quản lý có thể lưu snapshot, xuất báo cáo hoặc dùng kết quả để lập kế hoạch đặt hàng.
  \item Các nguyên liệu rủi ro cao, tiêu hao bất thường hoặc hao hụt lớn được làm nổi bật.
\end{itemize}

    \textbf{Luồng chính}
    \begin{enumerate}[leftmargin=1.8em]
  \item Quản lý chọn kỳ phân tích, chi nhánh và nhóm nguyên liệu cần xem.
  \item Hệ thống tải lịch sử tiêu hao, cảnh báo sắp hết hàng, dữ liệu hao hụt và thời gian cung ứng.
  \item Hệ thống tính mức tiêu hao bình quân, độ biến động theo ca và xác suất thiếu hàng trong kỳ kế tiếp.
  \item Hệ thống sinh lượng tồn mục tiêu, điểm đặt lại hàng và lượng nhập khuyến nghị.
  \item Quản lý rà soát các đề xuất, điều chỉnh nếu cần và lưu hoặc xuất báo cáo.
\end{enumerate}

    \textbf{Luồng thay thế}
    \begin{itemize}[leftmargin=1.8em]
  \item Nếu dữ liệu lịch sử chưa đủ dày, hệ thống chuyển sang quy tắc min-max và mức an toàn mặc định.
  \item Quản lý có thể so sánh nhiều kỳ như cuối tuần, ngày lễ hoặc giờ cao điểm để chọn phương án nhập hàng.
\end{itemize}

    \textbf{Ngoại lệ}
    \begin{itemize}[leftmargin=1.8em]
  \item Dữ liệu đơn vị tính không nhất quán hoặc thiếu ánh xạ nguyên liệu: hệ thống đánh dấu cần rà soát thủ công.
  \item Lịch sử giao dịch quá ít: hệ thống vẫn cho phép tạo đề xuất nhưng phải gắn cảnh báo mức tin cậy thấp.
\end{itemize}

    \textbf{Quy tắc nghiệp vụ}
    \begin{itemize}[leftmargin=1.8em]
  \item Đề xuất nhập hàng là khuyến nghị ra quyết định, không tự động sinh giao dịch mua hàng.
  \item Việc tính lượng nhập phải xét theo từng chi nhánh vì nhu cầu, lead time và mức hao hụt có thể khác nhau.
\end{itemize}
    \end{usecasebox}



    \begin{usecasebox}[title={UC-REP-01 --- Xem báo cáo doanh thu và vận hành}]
    \textbf{Tác nhân chính:} Quản lý\\
    \textbf{Mục tiêu:} Theo dõi doanh thu, giờ cao điểm, món bán chạy, độ trễ bếp và hiệu suất phục vụ.\\
    \textbf{Kích hoạt:} Quản lý mở bảng điều khiển phân tích hoặc yêu cầu xuất báo cáo.

    \medskip
    \textbf{Tiền điều kiện}
    \begin{itemize}[leftmargin=1.8em]
  \item Dữ liệu giao dịch, gọi món, bếp và tồn kho đã được tổng hợp vào mô hình đọc báo cáo.
  \item Người dùng có quyền xem báo cáo.
\end{itemize}

    \textbf{Hậu điều kiện}
    \begin{itemize}[leftmargin=1.8em]
  \item Quản lý có được ảnh chụp báo cáo hiện tại hoặc tệp xuất PDF/CSV.
  \item Các chỉ số dùng cho họp vận hành được chuẩn hóa.
\end{itemize}

    \textbf{Luồng chính}
    \begin{enumerate}[leftmargin=1.8em]
  \item Quản lý chọn khoảng thời gian, chi nhánh và bộ lọc báo cáo.
  \item Hệ thống truy xuất \texttt{ReportSnapshot} hoặc tính toán số liệu mới nhất.
  \item Hệ thống hiển thị doanh thu, số bàn quay vòng, món bán chạy, phiếu bếp chậm, tỷ lệ hủy món và hiệu suất nhân viên.
  \item Quản lý phân tích biểu đồ, đi sâu theo ca, trạm bếp hoặc khu vực.
  \item Nếu cần, quản lý chọn xuất báo cáo sang PDF/CSV.
\end{enumerate}

    \textbf{Luồng thay thế}
    \begin{itemize}[leftmargin=1.8em]
  \item Bảng điều khiển có thể hiển thị cảnh báo bất thường như thời gian ra món tăng mạnh hoặc hoàn tiền tăng đột biến.
  \item Quản lý có thể lưu bộ lọc ưa thích để dùng lại.
\end{itemize}

    \textbf{Ngoại lệ}
    \begin{itemize}[leftmargin=1.8em]
  \item Read model báo cáo chưa được làm mới: hệ thống hiển thị thời điểm đồng bộ cuối cùng.
  \item Khối lượng dữ liệu quá lớn: hệ thống chuyển truy vấn sang chế độ bất đồng bộ và thông báo khi hoàn tất.
\end{itemize}

    \textbf{Quy tắc nghiệp vụ}
    \begin{itemize}[leftmargin=1.8em]
  \item Báo cáo phân tích dùng dữ liệu tổng hợp, không truy cập trực tiếp vào luồng ghi giao dịch thời gian thực.
  \item Mỗi lần xuất báo cáo phải lưu dấu vết ai xuất và thời điểm xuất.
\end{itemize}
    \end{usecasebox}



    \begin{usecasebox}[title={UC-ADM-01 --- Quản lý menu và giá bán}]
    \textbf{Tác nhân chính:} Quản lý\\
    \textbf{Mục tiêu:} Tạo, cập nhật, tạm ẩn món, thay đổi giá và chương trình khuyến mãi.\\
    \textbf{Kích hoạt:} Quản lý mở màn hình cấu hình menu.

    \medskip
    \textbf{Tiền điều kiện}
    \begin{itemize}[leftmargin=1.8em]
  \item Người dùng có quyền quản trị menu.
  \item Chi nhánh mục tiêu đang được cấu hình và có danh mục món cơ sở.
\end{itemize}

    \textbf{Hậu điều kiện}
    \begin{itemize}[leftmargin=1.8em]
  \item Thực đơn mới hoặc thay đổi giá được lưu và có thể công bố theo lịch hiệu lực.
  \item Những món liên quan đến đơn gọi món đang mở vẫn giữ giá đã được chụp lại trước đó.
\end{itemize}

    \textbf{Luồng chính}
    \begin{enumerate}[leftmargin=1.8em]
  \item Quản lý chọn danh mục món hoặc món cụ thể.
  \item Hệ thống hiển thị thông tin giá, tình trạng bán, modifier và món liên quan.
  \item Quản lý chỉnh sửa tên, giá, tình trạng bán, hình thức khuyến mãi hoặc thuộc tính món.
  \item Quản lý lưu thay đổi và chọn công bố ngay hoặc theo lịch.
  \item Hệ thống kiểm tra tác động tới POS/KDS và phát hành bản menu mới.
\end{enumerate}

    \textbf{Luồng thay thế}
    \begin{itemize}[leftmargin=1.8em]
  \item Có thể công bố theo ca sáng, ca tối hoặc theo ngày lễ.
  \item Khuyến mãi có thể áp dụng theo món, danh mục hoặc mức hóa đơn.
\end{itemize}

    \textbf{Ngoại lệ}
    \begin{itemize}[leftmargin=1.8em]
  \item Dữ liệu cấu hình thiếu bắt buộc, ví dụ chưa gán trạm bếp hoặc thiếu thuế suất: hệ thống không cho công bố.
  \item Xung đột lịch hiệu lực giữa nhiều chương trình giá: hệ thống yêu cầu người dùng giải quyết.
\end{itemize}

    \textbf{Quy tắc nghiệp vụ}
    \begin{itemize}[leftmargin=1.8em]
  \item Không được thay đổi trực tiếp giá của đơn gọi món đã xác nhận.
  \item Mọi lần công bố thực đơn phải được ghi vào nhật ký kiểm toán và có phiên bản.
\end{itemize}
    \end{usecasebox}



    \begin{usecasebox}[title={UC-ADM-02 --- Quản lý vai trò và phân quyền}]
    \textbf{Tác nhân chính:} Quản trị viên, Quản lý\\
    \textbf{Mục tiêu:} Gán vai trò và quyền thao tác cho lễ tân, phục vụ, bếp, thu ngân và quản lý.\\
    \textbf{Kích hoạt:} Quản trị viên mở màn hình quản lý người dùng.

    \medskip
    \textbf{Tiền điều kiện}
    \begin{itemize}[leftmargin=1.8em]
  \item Tài khoản người dùng đã tồn tại hoặc đang được tạo mới.
  \item Chính sách RBAC và danh mục quyền đã được định nghĩa.
\end{itemize}

    \textbf{Hậu điều kiện}
    \begin{itemize}[leftmargin=1.8em]
  \item Tài khoản người dùng được gán đúng vai trò/quyền theo chi nhánh.
  \item Các thao tác nhạy cảm bị kiểm soát bởi \texttt{AuthorizationPolicy}.
\end{itemize}

    \textbf{Luồng chính}
    \begin{enumerate}[leftmargin=1.8em]
  \item Quản trị viên tìm người dùng cần cập nhật.
  \item Hệ thống hiển thị vai trò hiện tại, chi nhánh và phạm vi quyền.
  \item Quản trị viên thêm/bỏ vai trò hoặc gán quyền  theo chính sách.
  \item Hệ thống lưu thay đổi và yêu cầu người dùng đăng nhập lại nếu cần.
  \item Audit log ghi nhận toàn bộ thay đổi quyền.
\end{enumerate}

    \textbf{Luồng thay thế}
    \begin{itemize}[leftmargin=1.8em]
  \item Có thể gán vai trò tạm thời cho ca làm việc hoặc theo khoảng thời gian hiệu lực.
  \item Nhà hàng nhiều chi nhánh có thể áp dụng vai trò khác nhau theo từng địa điểm.
\end{itemize}

    \textbf{Ngoại lệ}
    \begin{itemize}[leftmargin=1.8em]
  \item Người dùng đang tự gỡ quyền quản trị cuối cùng của hệ thống: thao tác bị chặn.
  \item Vai trò xung đột với chính sách hiện hành: hệ thống yêu cầu xác nhận hoặc điều chỉnh.
\end{itemize}

    \textbf{Quy tắc nghiệp vụ}
    \begin{itemize}[leftmargin=1.8em]
  \item Nguyên tắc tối thiểu quyền được áp dụng mặc định.
  \item Hoàn tiền, ghi đè giá và điều chỉnh tồn kho là các quyền nhạy cảm phải được kiểm soát riêng.
\end{itemize}
    \end{usecasebox}



    \begin{usecasebox}[title={UC-ADM-03 --- Xem nhật ký kiểm toán và truy vết thao tác nhạy cảm}]
    \textbf{Tác nhân chính:} Quản lý, Quản trị viên\\
    \textbf{Mục tiêu:} Tra cứu, đối chiếu và truy vết các hành động nhạy cảm như hoàn tiền, hủy món, ghi đè giá, đổi quyền hoặc mở lại hóa đơn.\\
    \textbf{Kích hoạt:} Người có thẩm quyền mở màn hình kiểm toán để điều tra sự cố, rà soát cuối ca hoặc phục vụ kiểm tra nội bộ.

    \medskip
    \textbf{Tiền điều kiện}
    \begin{itemize}[leftmargin=1.8em]
  \item Hệ thống đang ghi nhận \texttt{AuditLog} cho các hành động nhạy cảm.
  \item Người dùng có quyền xem nhật ký kiểm toán ở phạm vi chi nhánh hoặc toàn hệ thống theo vai trò.
  \item Dữ liệu correlation id, thời gian và đối tượng bị tác động có thể truy xuất được.
\end{itemize}

    \textbf{Hậu điều kiện}
    \begin{itemize}[leftmargin=1.8em]
  \item Người dùng xem được chuỗi hành động liên quan đến một giao dịch hoặc một người thao tác.
  \item Hệ thống lưu lại dấu vết việc xem hoặc xuất audit log nếu chính sách yêu cầu.
  \item Có thể tạo báo cáo hoặc gắn cờ theo dõi cho bản ghi cần điều tra thêm.
\end{itemize}

    \textbf{Luồng chính}
    \begin{enumerate}[leftmargin=1.8em]
  \item Người dùng mở màn hình nhật ký kiểm toán.
  \item Hệ thống kiểm tra quyền truy cập vào dữ liệu kiểm toán nhạy cảm.
  \item Người dùng nhập bộ lọc như khoảng thời gian, người thao tác, thực thể, loại hành động hoặc mã liên kết.
  \item Hệ thống truy vấn các bản ghi audit bất biến và hiển thị danh sách kết quả phù hợp.
  \item Người dùng mở chi tiết một bản ghi để xem giá trị trước, giá trị sau, lý do, correlation id và ngữ cảnh liên quan.
  \item Nếu cần, người dùng xuất báo cáo hoặc đánh dấu follow-up để tiếp tục điều tra.
\end{enumerate}

    \textbf{Luồng thay thế}
    \begin{itemize}[leftmargin=1.8em]
  \item Người dùng có thể đi từ một hóa đơn, payment hoặc order cụ thể sang chuỗi audit liên quan.
  \item Các trường nhạy cảm như thông tin liên hệ hoặc dữ liệu thanh toán tham chiếu có thể được làm mờ theo quyền.
\end{itemize}

    \textbf{Ngoại lệ}
    \begin{itemize}[leftmargin=1.8em]
  \item Người dùng không đủ quyền xem chi tiết audit log: hệ thống từ chối truy cập hoặc chỉ cho xem bản rút gọn.
  \item Không có bản ghi phù hợp với bộ lọc: hệ thống thông báo không có kết quả và cho phép điều chỉnh bộ lọc.
  \item Kho audit tạm thời chậm hoặc không sẵn sàng: hệ thống chuyển sang chế độ chỉ đọc hạn chế và cảnh báo quản trị.
\end{itemize}

    \textbf{Quy tắc nghiệp vụ}
    \begin{itemize}[leftmargin=1.8em]
  \item Bản ghi kiểm toán đã phát sinh không được phép sửa nội dung theo cách làm mất dấu vết gốc.
  \item Việc xem hoặc xuất audit log mức nhạy cảm cao cũng phải được ghi nhận vào audit log.
\end{itemize}
    \end{usecasebox}



    \begin{usecasebox}[title={UC-OPS-01 --- Quản lý ca làm nhân viên}]
    \textbf{Tác nhân chính:} Quản lý\\
    \textbf{Mục tiêu:} Phân công nhân sự cho từng chi nhánh, khu vực và ca làm việc.\\
    \textbf{Kích hoạt:} Quản lý mở lịch làm việc và tạo/cập nhật ca.

    \medskip
    \textbf{Tiền điều kiện}
    \begin{itemize}[leftmargin=1.8em]
  \item Tài khoản nhân viên đã tồn tại.
  \item Mẫu ca làm và cấu hình khu vực phục vụ đã được tạo.
\end{itemize}

    \textbf{Hậu điều kiện}
    \begin{itemize}[leftmargin=1.8em]
  \item Một hoặc nhiều \texttt{ShiftAssignment} được tạo hoặc cập nhật.
  \item Nhân viên nhìn thấy ca được phân trong lịch làm việc.
\end{itemize}

    \textbf{Luồng chính}
    \begin{enumerate}[leftmargin=1.8em]
  \item Quản lý chọn ngày, chi nhánh và loại ca.
  \item Hệ thống hiển thị danh sách nhân viên khả dụng cùng kỹ năng/vai trò.
  \item Quản lý gán nhân viên cho ca và khu vực phụ trách.
  \item Hệ thống kiểm tra trùng ca, thiếu vị trí bắt buộc và giới hạn giờ làm.
  \item Quản lý lưu lịch; hệ thống gửi thông báo cho nhân viên.
\end{enumerate}

    \textbf{Luồng thay thế}
    \begin{itemize}[leftmargin=1.8em]
  \item Quản lý có thể sao chép lịch từ tuần trước rồi tinh chỉnh.
  \item Ca thiếu người có thể được đánh dấu chờ xác nhận hoặc mở đăng ký tự nguyện.
\end{itemize}

    \textbf{Ngoại lệ}
    \begin{itemize}[leftmargin=1.8em]
  \item Nhân viên bị phân ca trùng thời gian hoặc vượt giới hạn giờ làm: hệ thống chặn lưu.
  \item Người dùng không thuộc chi nhánh được chọn: hệ thống yêu cầu thay đổi phạm vi gán.
\end{itemize}

    \textbf{Quy tắc nghiệp vụ}
    \begin{itemize}[leftmargin=1.8em]
  \item Một nhân viên không được có hai ca chồng lấn ở cùng thời điểm.
  \item Ca bếp, ca thu ngân và ca quản lý phải đáp ứng số lượng tối thiểu theo chính sách vận hành.
\end{itemize}
    \end{usecasebox}

\section{Kiến trúc phần mềm}
